% !TeX program = xelatex
% UTF-8. Build twice with XeLaTeX. No shell escape or external assets required.
% Version 0.1. SPEC-GATE: PASS (2026-10-02). Baseline APPROVED by Yana Abramenko (2026-10-03).
\documentclass[11pt,a4paper]{article}
\usepackage[margin=22mm]{geometry}
\usepackage{fontspec}
\usepackage{polyglossia}
\setmainlanguage{ukrainian}
\setotherlanguage{english}
\IfFontExistsTF{Libertinus Serif}{
  \setmainfont{Libertinus Serif}
  \newfontfamily\cyrillicfont{Libertinus Serif}
}{
  \setmainfont{Times New Roman}
  \newfontfamily\cyrillicfont{Times New Roman}
}
\IfFontExistsTF{DejaVu Sans Mono}{\setmonofont{DejaVu Sans Mono}}{\setmonofont{Consolas}}
\usepackage{longtable,booktabs,array}
\usepackage{enumitem}
\usepackage{needspace}
\usepackage{fancyhdr}
\usepackage[unicode,hidelinks]{hyperref}
\hypersetup{pdftitle={Ветеринарна клініка — SRS 0.1},pdfauthor={Абраменко Яна},pdfsubject={SPEC-GATE: PASS; baseline затверджена людиною 03.10.2026}}
\setlength{\parindent}{0pt}
\setlength{\parskip}{5pt}
\setlength{\tabcolsep}{4pt}
\renewcommand{\arraystretch}{1.2}
\setlist[itemize]{leftmargin=*,itemsep=2pt,topsep=3pt}
\setlength{\emergencystretch}{3em}
\pagestyle{fancy}
\fancyhf{}
\fancyhead[L]{\small Ветеринарна клініка}
\fancyhead[R]{\small SRS 0.1}
\fancyfoot[L]{\small Затверджена baseline — SPEC-GATE: PASS}
\fancyfoot[R]{\thepage}
\setlength{\headheight}{15pt}
\renewcommand{\headrulewidth}{0.3pt}
\newcommand{\requirement}[4]{%
  \Needspace{7\baselineskip}%
  \subsubsection*{\hypertarget{#1}{\texttt{#1}}\quad #2}%
  #3\par
  \textbf{Критерій прийняття}\quad #4\par\medskip}
\newcommand{\scenario}[5]{%
  \Needspace{8\baselineskip}%
  \subsubsection*{\hypertarget{#1}{\texttt{#1}}\quad #2}%
  \textbf{Given}\quad #3\par
  \textbf{When}\quad #4\par
  \textbf{Then}\quad #5\par\medskip}
\begin{document}
\begin{titlepage}
\centering
{\Large Інформаційна система ветеринарної клініки\par}
\vspace{12mm}
{\LARGE\bfseries Специфікація вимог до програмного забезпечення\par}
\vspace{8mm}
{\large Software Requirements Specification\par}
\vspace{12mm}
{\large Лабораторна робота № 2\par}
{\large Проектування інформаційних систем\par}
\vspace{12mm}
Абраменко Яна\par
Версія 0.1\par
28 вересня 2026 року\par
\vfill
\textbf{Статус: затверджена базова специфікація}\par
SPEC-GATE: PASS (02.10.2026). Baseline затверджена Абраменко Яною 03.10.2026 для подальшого SDD-циклу.\par
\vspace{8mm}
\begin{minipage}{0.94\textwidth}
Документ формалізує рішення поетапного діалогу щодо однієї ветеринарної клініки.
Він містить межі, терміни, вимоги, критерії прийняття, сценарії Given–When–Then
і початкову матрицю простежуваності. Невирішені питання винесено до окремого
реєстру, а не замінено припущеннями coding agent.
\end{minipage}
\end{titlepage}
\tableofcontents
\clearpage
\section{Керування документом}
\begin{longtable}{>{\raggedright\arraybackslash}p{4cm} >{\raggedright\arraybackslash}p{11.6cm}}
\toprule
\textbf{Поле} & \textbf{Значення} \\ \midrule
\endfirsthead
\toprule
\textbf{Поле} & \textbf{Значення} \\ \midrule
\endhead
\bottomrule
\endfoot
Назва системи & Інформаційна система ветеринарної клініки \\
Версія і статус & 0.1; APPROVED; SPEC-GATE: PASS; baseline затверджена людиною 03.10.2026 \\
Підстава & ЛР1, погоджений SKED-діалог і завдання ЛР2 \\
Власниця рішень & Абраменко Яна \\
Мова документа й інтерфейсу & Українська \\
Методологія & SDD; обов’язковий TDD для функціональних змін \\
Цільове використання & Реальна мала клініка; також навчальна демонстрація \\
Файл базової специфікації & spec/SRS\_Vet\_Clinic\_v0.1.tex \\
Затвердження та коміт & Baseline затверджена Абраменко Яною 03.10.2026 після SPEC-GATE: PASS; коміт затвердження ще не створено \\
\end{longtable}

У цьому документі слова «повинна» та «повинен» визначають перевірювану поведінку чернетки. Усі вимоги походять із погоджених рішень; конкретизації, яких діалог не встановив, наведені в OPEN та DEP. Вони не є мовчазним дозволом агенту обрати суттєву поведінку. Структуру адаптовано за ULLANA; предметну область, стек і предметні обмеження прикладів не перенесено.

Доступний зразок ULLANA надано у PDF, а не у вихідному .tex. Тому використано стандартний синтаксис article, section, longtable, нумеровані вимоги й окремі розділи методології та приймання; точне відтворення авторських макросів зразка не заявляється.

\section{Призначення та початковий задум}
Система створює єдиний електронний простір однієї ветеринарної клініки. Клієнт бачить історію тварини, результати, щеплення, обробки, направлення, призначення й дієти та записується онлайн. Працівники ведуть електронні картки й розклад, отримують доступ до потрібної інформації без пошуку паперових карток. Зменшення телефонних звернень і відмова від паперових карток є цілями; автоматична міграція паперового архіву не погоджена. Цей документ призначений студентці, рецензенту, розробнику, тестувальнику й coding agent.

\section{Межі першої версії}
\begin{itemize}
\item Одна клініка без підтримки мережі філій; доступ через інтернет у браузері на ПК і смартфоні.
\item Клієнтський профіль незалежно від тварин; робочі облікові записи адміністратора й ветеринара, запрошення та відновлення доступу.
\item Картки тварин, основні дані, перевірка та пропозиції змін, смерть і припинення опіки зі збереженням історії.
\item Розклад на конкретні дати, двадцятихвилинні прийоми, онлайн-запис, адміністративне історичне внесення, статуси та їх контрольоване виправлення.
\item Медичні записи, результати та вкладення, процедури, текстові призначення й дієти, направлення, клінічні й домашні рекомендації.
\item Пошук клієнтів і тварин; обмежений організаційний доступ адміністратора; історія версій у погоджених випадках.
\item SMS-коди для трьох погоджених дій, технічна пошта про резервування, автоматичні копії та пробне відновлення.
\end{itemize}

\section{Поза межами поточної версії}
\begin{itemize}
\item Оплата послуг, фінансовий облік і складський облік препаратів.
\item Автоматичний обмін із лабораторіями чи іншими клініками; направлення й результати зберігаються без таких інтеграцій.
\item Окремі iOS/Android-застосунки, перемикання мов, мережа філій.
\item Передавання картки новому власнику, спільний доступ кількох поточних власників; порядок обслуговування зміни власника не визначається цією версією.
\item Повторюваний тижневий розклад, довідник послуг, довідники видів, порід і препаратів, автоматичні розрахунки лікування або ветеринарних графіків.
\item Завантаження сканів чи фото ветеринарного паспорта; зберігання копій або реквізитів документа клієнта при особистій перевірці.
\item Автоматичне надсилання запрошень; автоматичні нагадування клієнтам про процедури або візити не були погоджені й не включені.
\item Відмітка виконання домашньої процедури, окремий стан рекомендації «Виконана», чернетки й приховані від клієнта медичні записи.
\item Окремий список перевірки даних тварин; часткове схвалення пропозиції; кілька нерозглянутих пропозицій на тварину.
\item Редагування медичного запису не автором, зокрема після блокування автора; спеціальні коригувальні записи інших лікарів.
\item Загальне логування всіх бізнес-подій не вводиться. Погоджені історії виправлень зберігаються; логи розробки є окремим процесним артефактом.
\end{itemize}

\section{Джерела та пріоритет рішень}
\begin{longtable}{>{\raggedright\arraybackslash}p{2cm} >{\raggedright\arraybackslash}p{13.6cm}}
\toprule
\textbf{ID} & \textbf{Джерело і роль} \\ \midrule
\endfirsthead
\toprule
\textbf{ID} & \textbf{Джерело і роль} \\ \midrule
\endhead
\bottomrule
\endfoot
SRC-01 & ПІС\_Лаб1\_Абраменко.docx: задум, ролі, структура репозиторію, рішення щодо кількох тварин. Поточний стан приватного репозиторію в цій роботі не перевірявся. \\
SRC-02 & Повний текст завдання ЛР2, наданий у діалозі: SKED, SRS, AC, BDD, простежуваність, підсумок, SPEC-GATE, затверджувальний коміт. \\
SRC-03 & SRS\_Calorie\_Calculator\_v1.0.pdf: приклад структури та перевірюваності. Його виключення BDD до цього проєкту не застосовується. \\
SRC-04 & ullana\_specification.pdf, Version 0.1: приклад структури, SDD, TDD та простежуваності. JSONL-логування ULLANA не переноситься як функція клініки. \\
SRC-05 & Погоджений SKED-діалог 27–28.09.2026: уточнення предметної області, винятки, відхилені пропозиції та обмеження. Це основне джерело формулювань цієї чернетки. \\
\end{longtable}

Пізніше явне рішення людини в SRC-05 уточнює ранні домовленості. Зокрема, профіль може існувати без тварини, а адміністратор може створювати історичні прийоми поза розкладом. Інструкції й предметні вимоги прикладів не мають пріоритету над цим діалогом. Окремі розділи 3–5 Частини 1 з методологічного посібника не були надані; їх опрацювання не заявляється.

\section{Схема ідентифікаторів та домовленості}
\begin{longtable}{>{\raggedright\arraybackslash}p{4.4cm} >{\raggedright\arraybackslash}p{11.2cm}}
\toprule
\textbf{Префікс} & \textbf{Призначення} \\ \midrule
\endfirsthead
\toprule
\textbf{Префікс} & \textbf{Призначення} \\ \midrule
\endhead
\bottomrule
\endfoot
REQ-F-NNN & Функціональна вимога, включно з конкретною поведінкою зовнішніх інтерфейсів. \\
REQ-NF-NNN & Нефункціональна вимога або вимірюване обмеження якості. \\
REQ-DP-NNN & Development Process Requirement — вимога до процесу розробки. \\
AC-F/NF/DP-NNN & Критерій прийняття відповідної вимоги. \\
BDD-NNN & Сценарій Given–When–Then. Наявність сценарію не означає готового автоматизованого тесту. \\
OPEN-NNN / DEP-NNN & Питання для наступної перевірки / залежність реалізації чи експлуатації. \\
\end{longtable}

ID видані вперше в цій чернетці та залишаються стабільними надалі. Перестановка розділів не змінює ID; нові вимоги отримують нові номери. Вимога може мати комплексний критерій з кількома перевірками. У матриці посилання BDD означає релевантний приклад, а не доказ повного покриття всіх гілок AC. Позначка «—» у BDD означає пряму перевірку за AC, огляд або вимірювання, а не відсутність критерію.

\section{Глосарій}
\begin{longtable}{>{\raggedright\arraybackslash}p{4.3cm} >{\raggedright\arraybackslash}p{11.3cm}}
\toprule
\textbf{Термін} & \textbf{Визначення} \\ \midrule
\endfirsthead
\toprule
\textbf{Термін} & \textbf{Визначення} \\ \midrule
\endhead
\bottomrule
\endfoot
Клієнт & Особа, для якої існує профіль клініки; може ще не мати тварин чи доступу до входу. \\
Профіль клієнта & Ім’я, прізвище, телефон, необов’язкова пошта та зв’язки з тваринами; не тотожний сеансу чи паролю. \\
Обліковий запис & Засіб доступу зі своїми даними входу та станом. Особистий клієнтський і робочий доступи розділені. \\
Поточний власник & Клієнт, опіка якого над твариною не припинена; поточний власник не більше ніж один. \\
Колишній опікун & Клієнт із припиненою опікою, який зберігає погоджений перегляд картки та історії. \\
Картка тварини & Основні дані, зв’язки опіки, медична історія, рекомендації, направлення та відомості стану тварини. \\
Основні дані & Кличка, вид, стать, порода, народження, паспорт і єдиний номер мікрочипа; не медичні записи. \\
Підтверджені дані & Основні дані, прийняті працівником клініки; це не обов’язково документальна перевірка паспорта й не підтвердження всієї медичної історії. \\
Пропозиція змін & Один нерозглянутий набір нових значень підтверджених основних даних від клієнта. \\
Запис на прийом & Запис візиту з твариною, призначеним лікарем, початком, тривалістю й статусом; може бути внесений заднім числом. Це не обліковий запис користувача. \\
Підтверджений прийом & Збережений запис, для якого результат ще не зафіксовано; може мати минулий час. \\
Проведений прийом & Прийом, проведення якого окремо позначив призначений лікар за виконання вимог до опису. \\
Медичний запис & Відомості ветеринара про здоров’я, огляд, процедуру або результат із авторством та історією версій. \\
Помилковий медичний запис & Запис із явною позначкою автора й причиною, який не зараховується як коректний опис, але зберігається в історії. \\
Рекомендація & План подальшої дії «У клініці» чи «Вдома»; сама не бронює час і не має окремого стану «Виконана». \\
Рекомендація щодо запису & Окреме організаційне поле клінічної рекомендації, видиме адміністратору й клієнту. \\
Направлення & Інформація про аналіз, обстеження, спеціаліста чи установу; не створює рекомендації для запису автоматично. \\
Список очікування запису & Подання чинних клінічних рекомендацій, що потребують бронювання; не жива черга й не черга на вільний час. \\
Історія змін & Версії або записи виправлень усередині застосунку в погоджених випадках; не журнал роботи coding agent. \\
Логи розробки & Артефакти /logs з доказами тестування, роботи агента й рішень людини. \\
SDD / TDD / BDD & Розробка за специфікацією / розробка через попередні тести / опис поведінки через приклади Given–When–Then. \\
SPEC-GATE & Окрема контрольна перевірка специфікації з результатом PASS або FAIL. PASS передує окремому людському затвердженню baseline і не замінює його. \\
\end{longtable}

\section{Загальний опис і модель доступу}
Система є самостійним вебзастосунком однієї клініки. Тип бази даних, мови програмування, фреймворки, хостинг і постачальники інтеграцій не погоджені та не нав’язуються цим SRS. До 5 адміністраторів, 10 ветеринарів і 5000 клієнтів є погодженим орієнтиром; 50 активних користувачів — навантаження перевірки. Це не жорсткі обмеження кількості облікових записів.

\begin{longtable}{>{\raggedright\arraybackslash}p{3cm} >{\raggedright\arraybackslash}p{4.1cm} >{\raggedright\arraybackslash}p{4.1cm} >{\raggedright\arraybackslash}p{4.1cm}}
\toprule
\textbf{Дія} & \textbf{Клієнт} & \textbf{Адміністратор} & \textbf{Ветеринар} \\ \midrule
\endfirsthead
\toprule
\textbf{Дія} & \textbf{Клієнт} & \textbf{Адміністратор} & \textbf{Ветеринар} \\ \midrule
\endhead
\bottomrule
\endfoot
Основні дані тварини & До перевірки редагує; потім пропонує зміни; після смерті або припинення опіки лише переглядає & Створює, виправляє, підтверджує & Створює, виправляє, підтверджує \\
Медичні записи й файли & Перегляд доступних карток і версій & Не має медичного доступу & Переглядає всі; створює за правилами типу; редагує свої \\
Бронювання & Власні тварини, до початку & Клініка; історичні винятки & Не керує бронюваннями \\
Проведення прийому & Ні & Ні & Лише призначений лікар \\
Організаційна рекомендація & Перегляд і запис за своєю & Перегляд, список, запис & Створення; редагування автором \\
Смерть та опіка & Не встановлює & Фіксує й виправляє & Фіксує й виправляє \\
Робочі облікові записи & Ні & Створює, блокує, розблоковує з обмеженнями & Не адмініструє \\
\end{longtable}

Особа технічного обслуговування є зовнішньою експлуатаційною відповідальністю, а не четвертою роллю застосунку. Пошук не підвищує прав доступу. Колишній опікун зберігає доступ у межах погодженого правила; передавання іншому власнику не підтримується.

\subsection{Переходи статусів прийому}
\begin{longtable}{>{\raggedright\arraybackslash}p{5cm} >{\raggedright\arraybackslash}p{10.6cm}}
\toprule
\textbf{Перехід} & \textbf{Хто й умови} \\ \midrule
\endfirsthead
\toprule
\textbf{Перехід} & \textbf{Хто й умови} \\ \midrule
\endhead
\bottomrule
\endfoot
Створення → Підтверджено & Клієнт або адміністратор після успішного збереження; адміністратору дозволені історичні винятки. \\
Підтверджено → Проведений & Призначений лікар, не раніше початку, із коректним обов’язковим описом. \\
Підтверджено → Скасовано & Клієнт до початку; адміністратор до початку або після початку з причиною з боку клініки; також погоджене автоматичне скасування майбутніх записів. \\
Підтверджено → Неявка & Адміністратор після початку. \\
Скасовано / Неявка → Підтверджено & Адміністратор через виправлення, з причиною та перевірками конфліктів. \\
Проведений → Підтверджено & Призначений лікар через виправлення з причиною. \\
Перенесення & Статус залишається Підтверджено; лише до початку первісного прийому; стара бронь зберігається до успішного результату. \\
\end{longtable}

Прямий перехід між кінцевими статусами не замінює погодженого виправлення через «Підтверджено». Позначка помилкового медичного запису не є статусом прийому. Тривалість кожного інтервалу — 20 хвилин, включно з історичними прийомами.

\subsection{Умова включення рекомендації до списку}
Рекомендація включається, якщо одночасно: формат «У клініці»; її не скасовано; тварина має поточного власника та не має позначки смерті; за рекомендацією немає підтвердженого або проведеного прийому. Минулість рекомендованої дати не є підставою вилучення. Ця умова об’єднує погоджені правила; окремого стану «Виконана» не додає.

\section{Зовнішні інтерфейси}
\begin{longtable}{>{\raggedright\arraybackslash}p{4.2cm} >{\raggedright\arraybackslash}p{11.4cm}}
\toprule
\textbf{Інтерфейс} & \textbf{Погоджений контракт} \\ \midrule
\endfirsthead
\toprule
\textbf{Інтерфейс} & \textbf{Погоджений контракт} \\ \midrule
\endhead
\bottomrule
\endfoot
Браузерний & Український адаптивний інтерфейс; клієнтський кабінет, робочі подання, серверна авторизація. Піксельні макети не визначені. \\
SMS-сервіс & Український номер і код для реєстрації, відновлення або зміни номера; прийняття, відхилення чи тайм-аут до 10 секунд. Ключі постачальника не фіксуються у специфікації. \\
Технічна електронна пошта & Повідомлення про помилку резервного копіювання на адресу в налаштуваннях; до 5 хвилин після виявлення. \\
Вкладення & PDF/JPEG/PNG до погоджених меж; читання й завантаження під контролем доступу, включно з версіями. \\
Резервне сховище & Окреме від основного сервера, дані й файли, 7 діб, щогодинне копіювання та контрольне відновлення. \\
Ручне передавання запрошень & Система формує посилання; адміністратор передає його поза автоматичним каналом застосунку. \\
Зовнішні лабораторії й установи & Інтеграційного контракту немає; інформацію про направлення й результати вводить ветеринар. \\
\end{longtable}

\section{Функціональні вимоги}
\subsection{Клієнтський профіль і доступ}
\requirement{REQ-F-001}{Профіль клієнта без тварин}{Система повинна зберігати профіль клієнта незалежно від наявності тварин. Обов’язкові ім’я, прізвище та телефон; електронна пошта необов’язкова. Профіль може бути створений до надання доступу для входу.}{\hypertarget{AC-F-001}{\texttt{AC-F-001}}. Клієнт із заповненими обов’язковими полями та без пошти й тварин може мати профіль. Смерть або припинення опіки над останньою твариною не видаляє профіль.}

\requirement{REQ-F-002}{Формат та унікальність телефону}{Система повинна приймати лише +380 і дев’ять наступних цифр. Перед перевіркою прибираються пробіли, дужки та дефіси. Номер без коду країни не доповнюється автоматично. Один нормалізований номер належить лише одному профілю клієнта.}{\hypertarget{AC-F-002}{\texttt{AC-F-002}}. Значення +380 (67) 123-45-67 і +380671234567 визначають один номер; другий профіль не створюється. Номер 0671234567 та номер іншої країни відхиляються. Одночасні реєстрації не створюють дубля.}

\requirement{REQ-F-003}{Самостійна реєстрація}{Клієнт повинен мати можливість самостійно зареєструватися та встановити пароль після підтвердження володіння номером через SMS-код. За наявності профілю з таким номером новий профіль не створюється, а сам збіг номера не надає доступу до наявного профілю.}{\hypertarget{AC-F-003}{\texttt{AC-F-003}}. Правильний чинний код дозволяє завершити реєстрацію без тварини. Неправильний або прострочений код не завершує її. Існуючий профіль адміністратора не дублюється й не привласнюється через повторну реєстрацію.}

\requirement{REQ-F-004}{Профіль і запрошення від адміністратора}{Адміністратор повинен мати можливість створити профіль після телефонного або особистого звернення. Система створює одноразове посилання для встановлення пароля; адміністратор передає його особисто або вручну погодженим каналом. Строк дії 24 години; нове посилання скасовує попереднє. Завершення строку не видаляє профіль.}{\hypertarget{AC-F-004}{\texttt{AC-F-004}}. За чинним посиланням клієнт сам установлює пароль до вже створеного профілю. Повторне використання та використання після 24 годин відхиляються; нове посилання робить старе непридатним. Автоматичне надсилання запрошень відсутнє.}

\requirement{REQ-F-005}{Вхід клієнта}{Клієнт повинен входити за номером телефону та паролем. SMS-код під час звичайного входу не потрібний.}{\hypertarget{AC-F-005}{\texttt{AC-F-005}}. Правильні телефон і пароль відкривають клієнтський доступ без запиту SMS-коду; неправильні дані не надають доступу.}

\requirement{REQ-F-006}{Життєвий цикл SMS-коду}{Коди реєстрації, відновлення пароля та зміни номера повинні діяти 5 хвилин, допускати до 5 спроб введення та повторне надсилання не раніше ніж через 60 секунд. Код одноразовий і прив’язаний до конкретної дії. Новий код скасовує попередній для цієї дії.}{\hypertarget{AC-F-006}{\texttt{AC-F-006}}. До завершення 5 хвилин правильний код у межах дозволених спроб приймається; після завершення строку, вичерпання 5 невдалих спроб або успішного використання — ні. Код іншої дії й попередній замінений код відхиляються. Повтор через 59 секунд недоступний, через 60 — доступний.}

\requirement{REQ-F-007}{Самостійне відновлення пароля}{Клієнт повинен мати можливість установити новий пароль після перевірки SMS-коду на зареєстрованому номері.}{\hypertarget{AC-F-007}{\texttt{AC-F-007}}. Без правильного чинного коду пароль не змінюється. Після підтвердження й установлення нового пароля попередній пароль більше не забезпечує вхід.}

\requirement{REQ-F-008}{Зміна номера клієнтом}{Для зміни номера в кабінеті клієнт повинен увести поточний пароль і підтвердити новий вільний номер через SMS. До успішного завершення старий номер залишається чинним. Унікальність нового номера перевіряється під час збереження.}{\hypertarget{AC-F-008}{\texttt{AC-F-008}}. Неправильний пароль, непідтверджений або вже зайнятий новий номер залишають старий номер без змін. Після успіху вхід здійснюється з новим номером; конкурентна зміна іншого профілю не створює дубля.}

\requirement{REQ-F-009}{Особисте відновлення доступу}{За втрати кабінету й старого номера адміністратор повинен перевірити документ при особистому зверненні та явно зафіксувати результат. Зберігаються адміністратор, час і результат; копія та реквізити документа не зберігаються. Негативний результат не дозволяє почати зміну доступу. Після позитивного клієнт підтверджує новий вільний номер через SMS і сам установлює пароль за одноразовим посиланням. Посилання діє 24 години від створення. Посилання, SMS-код і результат підтвердження прив’язані до конкретного профілю, нового номера та цієї операції відновлення. Успішне підтвердження номера чинне до завершення строку посилання, без окремого таймера. Нове посилання скасовує попередню незавершену операцію та потребує нового підтвердження номера; повторне надсилання SMS замінює лише код за REQ-F-006. Незавершена або прострочена процедура не змінює старі дані доступу. Перед остаточним спільним збереженням повторно перевіряється вільність нового номера.}{\hypertarget{AC-F-009}{\texttt{AC-F-009}}. До завершення обох кроків номер і пароль не змінюються. Після обох кроків дані доступу змінюються разом. Адміністратор не задає пароль. У збережених даних перевірки немає копії чи реквізитів документа. Процедура достатності доказів визначена залежністю DEP-004. За посилання, створеного о 09:00, підтвердження номера о 10:00 дозволяє встановити пароль о 10:10 того самого дня. Після завершення 24 годин або успішного використання посилання завершити операцію за ним не можна. Код чи підтвердження іншого профілю, номера або операції не приймаються. Після перевидачі посилання старе підтвердження не діє, потрібне нове. Якщо номер зайнято до остаточного збереження або процедура прострочена, старі номер і пароль залишаються без змін.}

\requirement{REQ-F-010}{Завершення сеансів при зміні доступу}{Після успішного відновлення пароля клієнта або працівника, а також зміни номера клієнта система повинна завершити всі попередні сеанси. Потрібний новий вхід.}{\hypertarget{AC-F-010}{\texttt{AC-F-010}}. Сеанси на двох пристроях після успішної операції не дозволяють наступного захищеного запиту; після повторного входу нові дані працюють. Невдала операція не змінює чинних даних доступу.}

\requirement{REQ-F-011}{Редагування профілю клієнта}{Клієнт повинен мати можливість змінювати власні ім’я, прізвище та пошту без підтвердження працівником. Адміністратор може виправляти ці поля на прохання клієнта. Ветеринар не може їх редагувати. Зміна телефону виконується лише за окремими процедурами REQ-F-008 і REQ-F-009.}{\hypertarget{AC-F-011}{\texttt{AC-F-011}}. Дозволені зміни зберігаються клієнтом або адміністратором; запит ветеринара відхиляється. Звичайне редагування профілю не обходить SMS-перевірку номера.}

\subsection{Облікові записи працівників}
\requirement{REQ-F-012}{Створення працівників}{Перший адміністратор повинен створюватися при початковому налаштуванні. Наступних адміністраторів і ветеринарів створює адміністратор через систему. Відкрита реєстрація працівників відсутня.}{\hypertarget{AC-F-012}{\texttt{AC-F-012}}. Після початкового налаштування перший адміністратор може створити працівника; клієнт і неавторизований користувач не можуть створити робочий обліковий запис.}

\requirement{REQ-F-013}{Логін та вхід працівника}{Працівник повинен входити за окремим унікальним логіном і паролем. Логін має 3–50 символів: латинські літери, цифри, крапка, дефіс, підкреслення; без пробілів. Регістр не впливає на вхід та унікальність.}{\hypertarget{AC-F-013}{\texttt{AC-F-013}}. Логіни Vet.Ivan і vet.ivan не можуть належати різним робочим обліковим записам. Довжини 3 і 50 дозволені, 2 і 51 — ні. Пробіл і кирилиця в логіні відхиляються.}

\requirement{REQ-F-014}{Пароль працівника через посилання}{Для первинного пароля та його відновлення адміністратор повинен створювати одноразове посилання на 24 години; нове скасовує попереднє. Працівник установлює пароль сам. При втраті пароля він звертається до адміністратора.}{\hypertarget{AC-F-014}{\texttt{AC-F-014}}. Адміністратор може видати посилання, але не задає пароль. Повторне, прострочене та замінене посилання відхиляються. Відновлення завершує попередні сеанси. Аварійний випадок єдиного адміністратора позначено OPEN-003.}

\requirement{REQ-F-015}{Блокування працівника}{Адміністратор повинен мати можливість блокувати працівника із завершенням активних сеансів. Його записи й авторство зберігаються. Заблокований ветеринар недоступний для нових майбутніх записів і перенесення до нього. Наявні майбутні записи не скасовуються автоматично та опрацьовуються адміністратором окремо.}{\hypertarget{AC-F-015}{\texttt{AC-F-015}}. Після блокування захищені запити працівника відхиляються, його немає у виборі лікарів для майбутнього запису; наявні прийоми та медичні записи зберігаються. Історичний виняток перевіряється за REQ-F-053.}

\requirement{REQ-F-016}{Розблокування працівника}{Адміністратор повинен мати можливість розблокувати працівника. Права його ролі стають доступними після нового входу; завершені сеанси не відновлюються.}{\hypertarget{AC-F-016}{\texttt{AC-F-016}}. Старий сеанс після розблокування не працює, новий успішний вхід надає права незмінної ролі.}

\requirement{REQ-F-017}{Захист адміністративного доступу}{Система повинна забороняти самоблокування адміністратора та блокування останнього незаблокованого адміністратора, який уже встановив пароль.}{\hypertarget{AC-F-017}{\texttt{AC-F-017}}. Запит самоблокування відхиляється. Обліковий запис, який ще не встановив пароль, не дозволяє заблокувати останнього працездатного адміністратора. Паралельні дії не залишають систему без такого адміністратора.}

\requirement{REQ-F-018}{Розділення облікових записів і ролей}{Робочий обліковий запис повинен мати одну незмінну роль: адміністратор або ветеринар. Для іншої ролі створюється окремий обліковий запис. Особистий клієнтський доступ працівника є окремим; перемикання ролей усередині облікового запису відсутнє.}{\hypertarget{AC-F-018}{\texttt{AC-F-018}}. Працівник-клієнт має окремі входи; жоден із них не надає сукупності прав двох ролей. Зміна ролі наявного робочого облікового запису недоступна.}

\subsection{Картки тварин та перевірка даних}
\requirement{REQ-F-019}{Основні дані тварини}{Картка повинна містити обов’язкові кличку, текстовий вид та стать зі значень «Самець», «Самка», «Невідомо». Порода є текстовою або явно невідомою. Народження задається точною датою, приблизним роком або «невідомо». Точна дата народження не може бути пізнішою за поточну дату Europe/Kyiv, приблизний рік — пізнішим за поточний рік. Додаткової межі давності не встановлено. Номер паспорта необов’язковий. Фото чи скан паспорта не додаються.}{\hypertarget{AC-F-019}{\texttt{AC-F-019}}. Відсутність клички, виду або вибору статі не дозволяє зберегти картку. Невідомі порода й народження та лише рік народження приймаються без підстановки дня й місяця. Довідників видів і порід немає. Поточна дата та поточний рік народження приймаються, майбутня дата й наступний рік — відхиляються за часом клініки; невідоме народження залишається допустимим.}

\requirement{REQ-F-020}{Номер мікрочипа}{У картці повинно бути одне необов’язкове поле мікрочипа: рівно 15 цифр зі збереженням початкових нулів. Значення унікальне серед карток, включно з картками померлих тварин і припиненої опіки. Дубль блокується також при схваленні зміни; клієнту пропонується звернутися до адміністратора без розкриття чужих даних.}{\hypertarget{AC-F-020}{\texttt{AC-F-020}}. Значення 001234567890123 зберігає нулі. 14 або 16 цифр і літери відхиляються. З двох одночасних спроб зберегти один номер успішна лише одна. Друге повідомлення не містить даних іншої картки.}

\requirement{REQ-F-021}{Зв’язок тварини з клієнтом}{Клієнт може мати нуль, одну або кілька тварин. Тварина повинна мати не більше одного поточного власника; після припинення опіки поточний власник може бути відсутній. Колишній опікун зберігає погоджений перегляд історії. Передавання картки та спільний доступ кількох поточних власників не підтримуються.}{\hypertarget{AC-F-021}{\texttt{AC-F-021}}. До клієнта додаються дві тварини; відсутність тварин не перешкоджає профілю. Після припинення опіки поточний власник відсутній, але колишній опікун бачить картку. Автоматичного передавання чи створення дубля немає.}

\requirement{REQ-F-022}{Самостійне додавання тварини}{Клієнт повинен мати можливість додати тварину після явного підтвердження наявності ветеринарного паспорта; номер паспорта вводиться, якщо відомий. Дані залишаються непідтвердженими до перевірки працівником під час першого звернення. До першого підтвердження клієнт може самостійно редагувати основні дані, крім випадків смерті або припинення опіки за REQ-F-031.}{\hypertarget{AC-F-022}{\texttt{AC-F-022}}. Без підтвердження наявності паспорта самостійне додавання не завершується. За підтвердженої наявності та невідомого номера картка додається як непідтверджена, видима клієнту та доступна для його редагування.}

\requirement{REQ-F-023}{Додавання тварини працівником}{Адміністратор або ветеринар повинен мати можливість під час особистого звернення створити картку зі слів власника, навіть без паспорта. Внесені основні дані одразу підтверджені: підтвердження означає їх прийняття працівником.}{\hypertarget{AC-F-023}{\texttt{AC-F-023}}. Обидві ролі створюють картку без номера й наявності паспорта; дані мають позначку підтвердження. Обов’язкові поля REQ-F-019 зберігають силу.}

\requirement{REQ-F-024}{Первинна перевірка даних}{Адміністратор або ветеринар повинен мати можливість виправити дані перед першим підтвердженням. Якщо інформації недостатньо, дані залишаються непідтвердженими з поясненням для клієнта; окремого статусу відхиленої картки немає.}{\hypertarget{AC-F-024}{\texttt{AC-F-024}}. Після перевірки виправлений варіант стає підтвердженим. При недостатніх даних пояснення видно клієнту, який може виправити їх самостійно.}

\requirement{REQ-F-025}{Робота до підтвердження тварини}{Непідтверджені основні дані не повинні блокувати перший онлайн-запис, створення медичних записів або позначення прийому проведеним за решти виконаних умов.}{\hypertarget{AC-F-025}{\texttt{AC-F-025}}. Для непідтвердженої тварини створюється прийом, лікар зберігає причину звернення й результати огляду та окремо позначає його проведеним.}

\requirement{REQ-F-026}{Пропозиція зміни підтверджених даних}{Після підтвердження клієнт повинен подавати пропозицію змін замість прямої заміни даних. Допускається лише одна нерозглянута пропозиція на тварину. Чинні дані залишаються в картці, пропозиція показується окремо клієнту та працівникам перевірки; у картці є позначка потреби перевірки. Окремого спільного списку немає.}{\hypertarget{AC-F-026}{\texttt{AC-F-026}}. Подання пропозиції не змінює чинних значень. Друга пропозиція до рішення недоступна. Клієнт, адміністратор і ветеринар бачать відмінність між чинними та запропонованими даними.}

\requirement{REQ-F-027}{Рішення щодо пропозиції}{Адміністратор або ветеринар повинен схвалювати чи відхиляти всю пропозицію цілком; дистанційна перевірка дозволена. Причина відхилення обов’язкова й видима клієнту. Схвалення замінює чинні поля, відхилення залишає їх без змін. Часткове схвалення не підтримується.}{\hypertarget{AC-F-027}{\texttt{AC-F-027}}. Після відхилення з причиною клієнт може подати нову пропозицію. Відхилення без причини не зберігається. Схвалення змінює всі запропоновані поля лише після валідації, включно з мікрочипом.}

\requirement{REQ-F-028}{Безпосереднє виправлення працівником}{Адміністратор або ветеринар повинен мати можливість самостійно виправити підтверджені основні дані. Якщо хоча б одне змінене поле є у нерозглянутій пропозиції клієнта, вся пропозиція відхиляється як неактуальна з причиною. Зміна іншого поля не закриває пропозицію.}{\hypertarget{AC-F-028}{\texttt{AC-F-028}}. За пропозиції клички та породи виправлення породи працівником відхиляє всю пропозицію з видимою причиною. Виправлення лише номера паспорта залишає таку пропозицію відкритою.}

\requirement{REQ-F-029}{Одночасний розгляд пропозиції}{Система повинна приймати лише перше успішно збережене рішення щодо нерозглянутої пропозиції. Інший працівник отримує повідомлення про вже виконаний розгляд та актуальний стан.}{\hypertarget{AC-F-029}{\texttt{AC-F-029}}. При одночасних схваленні й відхиленні зберігається рівно одне рішення; другий запит не перезаписує його.}

\subsection{Смерть і припинення опіки}
\requirement{REQ-F-030}{Смерть і припинення опіки}{Адміністратор або ветеринар повинен окремо фіксувати смерть і припинення опіки; клієнт сам їх не встановлює. Це незалежні відомості. Зберігаються дата й час події або лише дата, якщо час невідомий, та окремо момент внесення. Дата й час події не можуть бути пізнішими за поточний момент Europe/Kyiv; якщо час невідомий, дата не може бути пізнішою за поточну дату. Додаткової межі давності не встановлено. Позначка діє одразу після збереження.}{\hypertarget{AC-F-030}{\texttt{AC-F-030}}. За невідомого часу система зберігає лише дату без вигаданого часу. Клієнтський запит установлення позначки відхиляється. Установлення однієї позначки не підміняє другу. Майбутні дата/час події, а за невідомого часу — майбутня дата, відхиляються. Подія на поточний момент або лише з поточною датою приймається.}

\requirement{REQ-F-031}{Наслідки смерті та припинення опіки}{Картка, медична історія та профіль клієнта повинні зберігатися. Клієнт зберігає перегляд, але не створює нових прийомів і пропозицій змін. Основні дані доступні йому лише для перегляду незалежно від підтвердження; пряме редагування непідтверджених даних також заборонене. Картки відображаються в окремому розділі кабінету. Нерозглянута пропозиція автоматично відхиляється з відповідною причиною.}{\hypertarget{AC-F-031}{\texttt{AC-F-031}}. Після кожної події картка доступна в окремому розділі; спроби клієнта забронювати чи запропонувати зміни відхиляються. Чинна пропозиція відхилена з причиною; історію не видалено. Для непідтвердженої картки після кожної події пряме редагування основних даних клієнтом також відхиляється.}

\requirement{REQ-F-032}{Автоматичне скасування майбутніх прийомів}{Після внесення смерті або припинення опіки система повинна автоматично скасувати лише підтверджені прийоми, які на момент внесення ще не почалися, з відповідною видимою клієнту причиною. Прийоми, що вже почалися або відбулися, автоматично не змінюються. Дата й час самої події окремо використовуються для перевірки історичних прийомів за REQ-F-052. Рекомендації зберігаються в історії без показу у списку очікування.}{\hypertarget{AC-F-032}{\texttt{AC-F-032}}. При внесенні о 10:00 події, що сталася о 09:00, підтверджений запис на 11:00 скасовується, на 10:00 або 09:00 — ні; проведений прийом не змінюється. Якщо створення майбутнього запису випередило внесення позначки, він скасовується; якщо позначка збережена першою, створення блокується. Рекомендація не лишається у списку.}

\requirement{REQ-F-033}{Виправлення позначок стану тварини}{Адміністратор або ветеринар повинен мати можливість виправити помилкову позначку з обов’язковою причиною та історією зміни. Інша незалежна позначка не змінюється. Запис і пропозиції змін знову дозволені лише за поточного власника та відсутності смерті. Скасовані прийоми не відновлюються автоматично.}{\hypertarget{AC-F-033}{\texttt{AC-F-033}}. Після виправлення смерті при чинному припиненні опіки запис усе ще заборонений. Коли обидві умови доступу виконані, чинні клінічні рекомендації без підтвердженого або проведеного прийому повертаються до списку навіть із минулою датою. Скасовані рекомендації не повертаються.}

\subsection{Пошук}
\requirement{REQ-F-034}{Пошук клієнтів}{Адміністратори й ветеринари повинні шукати клієнтів за повним нормалізованим телефоном або частиною імені чи прізвища без урахування регістру. Результати містять ім’я, прізвище, телефон; профіль відкриває перелік пов’язаних тварин.}{\hypertarget{AC-F-034}{\texttt{AC-F-034}}. Частина прізвища знаходить відповідні профілі незалежно від регістру; частина телефону не є точним збігом. Однойменних клієнтів можна розрізнити за телефоном. Знайдений профіль не розширює прав ролі.}

\requirement{REQ-F-035}{Пошук тварин}{Працівники повинні шукати тварин за частиною клички без урахування регістру або точним повним мікрочипом. Результати містять кличку, вид і поточного власника; якщо його немає — «Поточний власник відсутній». Померлі тварини й припинена опіка включаються з позначками.}{\hypertarget{AC-F-035}{\texttt{AC-F-035}}. Запит «мур» знаходить «Мурка» та інші відповідні клички; унікальний мікрочип знаходить одну картку. Картка припиненої опіки не показує колишнього опікуна як поточного власника.}

\subsection{Розклад і бронювання}
\requirement{REQ-F-036}{Розклад на конкретні дати}{Лише адміністратор повинен задавати робочі інтервали лікарів на конкретні дати. Кілька інтервалів можуть визначати перерву між ними. Повторюваного тижневого розкладу немає.}{\hypertarget{AC-F-036}{\texttt{AC-F-036}}. Адміністратор задає 09:00–13:00 та 14:00–18:00; у перерві час для звичайного запису не пропонується. Ветеринар і клієнт не можуть змінювати розклад.}

\requirement{REQ-F-037}{Тривалість і доступні інтервали}{Звичайні записи повинні формуватися послідовно по 20 хвилин від початку кожного робочого інтервалу. Прийом повністю вміщується в нього; залишок менше 20 хвилин не пропонується. Суміжні інтервали не перетинаються.}{\hypertarget{AC-F-037}{\texttt{AC-F-037}}. Для 09:00–09:50 доступні початки 09:00 і 09:20, але не 09:40. Прийоми 10:00–10:20 та 10:20–10:40 не конфліктують.}

\requirement{REQ-F-038}{Зміна розкладу з наявними записами}{При зміні розкладу система повинна перевіряти лише підтверджені прийоми, які ще не почалися. Вони повинні повністю вміщуватися в нові робочі інтервали, але не зобов’язані збігатися з новою двадцятихвилинною сіткою; їхній час залишається незмінним. Якщо прийом не вміщується, зміна блокується з показом запису-перешкоди до його окремого перенесення або скасування. Нові записи формуються за новою сіткою з перевіркою перетинів. Історичні прийоми зміну розкладу не блокують.}{\hypertarget{AC-F-038}{\texttt{AC-F-038}}. За підтвердженого майбутнього прийому 10:00–10:20 вилучення цього робочого часу не зберігається й показує прийом. Після належного перенесення або скасування зміна зберігається. За прийому 10:00–10:20 зміна інтервалу 09:00–13:00 на 09:10–13:00 дозволена: прийом зберігає час 10:00, а нові слоти з перетином недоступні. Історичний прийом поза розкладом не блокує зміну.}

\requirement{REQ-F-039}{Створення звичайного запису}{Клієнт повинен обирати свою тварину, конкретного доступного ветеринара та вільний час. Адміністратор може створювати записи клієнтів. Після успішного збереження запис одразу має статус «Підтверджено», без окремого погодження. Довідник послуг не потрібний.}{\hypertarget{AC-F-039}{\texttt{AC-F-039}}. Доступний інтервал успішно зберігається й одразу показується підтвердженим. Для клієнта не потрібна дія адміністратора або вибір послуги.}

\requirement{REQ-F-040}{Майбутній час клієнтського запису}{Клієнт повинен створювати запис і вибирати новий час при перенесенні лише до настання цього часу. Адміністративне внесення минулих прийомів визначене окремо.}{\hypertarget{AC-F-040}{\texttt{AC-F-040}}. О 10:00 клієнт не може створити запис із початком 10:00 або раніше; доступний 10:20 дозволений за решти умов. Серверна перевірка відхиляє застарілий вибір із відкритої форми.}

\requirement{REQ-F-041}{Перетини та конкурентне бронювання}{Система повинна забороняти часові перетини прийомів одного ветеринара та підтверджених або проведених прийомів однієї тварини незалежно від лікаря, включно зі створенням історичного прийому. Скасовані записи не займають інтервал. Позначення «Неявка» не повинно автоматично звільняти часовий інтервал для інших історичних прийомів. Перевірка конфліктів із підтвердженими та проведеними прийомами зберігається. При одночасному виборі часу зберігається лише перше успішне бронювання; іншим повідомляється про зайнятість. Для екстрених випадків винятку немає.}{\hypertarget{AC-F-041}{\texttt{AC-F-041}}. Два конкурентні запити на один інтервал лікаря дають один запис. Одній тварині не створюється історичний прийом, що перетинає її підтверджений або проведений прийом у будь-якого лікаря. За неявки до лікаря на 10:00–10:20 інший історичний прийом до нього на цей інтервал не створюється. Суміжні інтервали дозволені. Відновлення минулих записів враховує також проведені прийоми за REQ-F-050.}

\requirement{REQ-F-042}{Захист від повторного надсилання}{Повтор тієї самої спроби створення запису після подвійного натискання чи обриву зв’язку повинен повертати вже створений запис, а не створювати дубль або повідомляти про зайнятість власним успішним записом.}{\hypertarget{AC-F-042}{\texttt{AC-F-042}}. Після успішного збереження з втраченою відповіддю повтор тієї самої спроби повертає той самий запис; у системі залишається один запис.}

\requirement{REQ-F-043}{Перенесення запису}{Клієнт повинен переносити власні записи, адміністратор — записи клініки, лише до початку первісного прийому та лише на майбутній час. Історичне внесення залишається окремою операцією. Новий час бронюється успішно до звільнення старого; невдача залишає старий запис без змін. Зв’язок із рекомендацією зберігається. Скасування рекомендації не перешкоджає перенесенню вже створеного прийому; рекомендація залишається скасованою. Після початку перенесення заборонене обом ролям.}{\hypertarget{AC-F-043}{\texttt{AC-F-043}}. За зайнятого нового часу старий прийом і його зв’язок залишаються. За успіху звільняється старий час і рекомендація не повертається до списку. У момент початку первісного прийому перенесення відхиляється. Обидві ролі не можуть перенести майбутній прийом на минулий час. Перенесення на майбутній доступний час за скасованої рекомендації дозволене, зв’язок і скасування рекомендації зберігаються.}

\requirement{REQ-F-044}{Перенесення до іншого ветеринара}{Клієнт для запису своєї тварини та адміністратор для записів клініки повинні мати можливість перенести майбутній прийом до іншого доступного ветеринара лише за відсутності пов’язаного медичного запису. Зміна лікаря не змінює авторство медичних даних. Заблокований ветеринар не є доступним цільовим лікарем.}{\hypertarget{AC-F-044}{\texttt{AC-F-044}}. Клієнт для своєї тварини й адміністратор можуть атомарно перенести прийом без медичного запису до іншого доступного лікаря на майбутній вільний час. Невдала спроба залишає попередній запис без змін. За наявності медичного запису зміна лікаря блокується. Зв’язок із рекомендацією зберігається.}

\requirement{REQ-F-045}{Скасування до початку}{Клієнт повинен мати можливість скасувати запис своєї тварини до початку, адміністратор — будь-який запис до початку. Лікар бронюваннями не керує. Скасування звільняє час і запускає правила повернення рекомендацій.}{\hypertarget{AC-F-045}{\texttt{AC-F-045}}. Клієнт скасовує власний майбутній запис, але не чужий і не запис, початок якого настав. Ветеринар не може виконати цю дію. Після скасування час доступний за чинного розкладу.}

\subsection{Результат прийому та виправлення}
\requirement{REQ-F-046}{Позначення проведеного прийому}{Лише призначений ветеринар повинен позначати прийом «Проведений», не раніше його початку та після збереження коректного опису з причиною звернення і результатами огляду. Збереження медичного запису саме по собі статус не змінює.}{\hypertarget{AC-F-046}{\texttt{AC-F-046}}. Без обов’язкового опису, до початку або від іншого лікаря чи адміністратора дія відхиляється. Після виконання умов окрема дія призначеного лікаря змінює статус.}

\requirement{REQ-F-047}{Неявка}{Адміністратор повинен мати можливість установити «Неявка» після початку запланованого прийому. Пов’язана рекомендація повертається до списку очікування з урахуванням винятків щодо стану тварини та скасованих рекомендацій.}{\hypertarget{AC-F-047}{\texttt{AC-F-047}}. До початку неявка не встановлюється; після початку адміністратор установлює її й відповідна чинна рекомендація знову доступна для запису. Пряме переписування кінцевого статусу не обходить правила виправлення.}

\requirement{REQ-F-048}{Скасування після початку з боку клініки}{Адміністратор повинен мати можливість скасувати прийом після початку через проблему клініки, якщо він ще не проведений, не скасований і не позначений неявкою. Причина обов’язкова й видима клієнту. Пов’язана рекомендація обробляється за загальними правилами скасування.}{\hypertarget{AC-F-048}{\texttt{AC-F-048}}. Підтверджений прийом після початку скасовується з причиною; без причини або за кінцевого статусу дія відхиляється. Перенесення після початку залишається недоступним.}

\requirement{REQ-F-049}{Виправлення статусу прийому}{Адміністратор повинен виправляти власні та чужі «Скасовано» й «Неявка» до «Підтверджено». Лише призначений лікар виправляє помилковий «Проведений» до «Підтверджено». Зберігаються попередній і новий статус, автор, час та обов’язкова причина. Адміністратор не може встановити «Проведений».}{\hypertarget{AC-F-049}{\texttt{AC-F-049}}. Виправлення без причини або від неналежної ролі відхиляється. Дозволене виправлення має всі п’ять реквізитів історії. Після виправлення неявки факт проведення окремо встановлює лікар.}

\requirement{REQ-F-050}{Перевірки відновлення запису}{Відновлення прийому не обходить обмежень щодо стану тварини та доступності ветеринара. Відновлення майбутнього запису повинно вимагати вільного первісного інтервалу, відповідності чинному розкладу, поточного власника, відсутності позначки смерті та незаблокованого ветеринара. Для минулого запису застосовуються погоджені правила історичного внесення REQ-F-052 і REQ-F-053; поточний розклад не перевіряється; перевіряються конфлікти з підтвердженими та проведеними прийомами й обмеження рекомендації. За іншого підтвердженого або проведеного прийому за тією самою рекомендацією відновлення блокується, без автоматичного скасування іншого запису.}{\hypertarget{AC-F-050}{\texttt{AC-F-050}}. Майбутній запис у вже зайнятий або вилучений із розкладу час не відновлюється. Минулий неконфліктний запис відновлюється навіть без такого часу в поточному розкладі. Конфліктний запис або другий активний зв’язок із рекомендацією не створюється. Майбутній запис не відновлюється після смерті тварини, припинення опіки або блокування лікаря. Минулий запис заблокованого лікаря може бути відновлений за допустимої хронології REQ-F-052 та решти виконаних перевірок; за недопустимої хронології відновлення відхиляється.}

\subsection{Історичне внесення прийомів}
\requirement{REQ-F-051}{Адміністративне внесення історичного прийому}{Адміністратор повинен мати можливість створювати прийом на фактичний час, що вже настав, поза розкладом та без обмеження давності. Тривалість залишається 20 хвилин, статус — «Підтверджено». Перетини заборонені; автоматичного «Проведений» немає.}{\hypertarget{AC-F-051}{\texttt{AC-F-051}}. Неконфліктний минулий прийом поза робочим інтервалом зберігається на 20 хвилин як підтверджений. Той самий запит клієнта відхиляється. Екстрена причина не дозволяє перетину.}

\requirement{REQ-F-052}{Історичний прийом після смерті або припинення опіки}{Адміністратор повинен мати можливість внести прийом, що передував події. Для припинення опіки всі 20 хвилин завершуються не пізніше події; для смерті дозволений також прийом, під час якого вона настала. Якщо час події невідомий, прийом того самого дня потребує явного підтвердження адміністратора з автором, часом і причиною допустимості.}{\hypertarget{AC-F-052}{\texttt{AC-F-052}}. За припинення опіки о 10:20 прийом 10:00–10:20 дозволений, 10:10–10:30 — ні. За смерті о 10:10 прийом 10:00–10:20 дозволений. Для події лише з датою прийом того самого дня не зберігається без зафіксованого підтвердження.}

\requirement{REQ-F-053}{Історичний прийом заблокованого лікаря}{Адміністратор повинен мати можливість створити минулий прийом заблокованого ветеринара. Він залишається «Підтвердженим» до можливого розблокування призначеного лікаря. Інший ветеринар не отримує права описати цей прийом або позначити його проведеним.}{\hypertarget{AC-F-053}{\texttt{AC-F-053}}. Історичний запис створюється, але новий майбутній запис до того самого заблокованого лікаря відхиляється. Інший лікар не може завершити історичний прийом.}

\subsection{Медична інформація}
\requirement{REQ-F-054}{Доступ до медичної частини}{Усі ветеринари повинні переглядати всі медичні картки. Клієнт із доступом до картки бачить медичні записи одразу після збереження та їхні попередні версії й причини виправлень. Чернеток і прихованих від клієнта записів немає. Адміністратор не бачить діагнозів, результатів, дієт, схем лікування, медичних вкладень та їхніх версій.}{\hypertarget{AC-F-054}{\texttt{AC-F-054}}. Збережений запис доступний належному клієнту й іншому ветеринару, але не сторонньому клієнту або адміністратору. Попередні версії підлягають тим самим обмеженням. Адміністратор бачить лише погоджену організаційну частину REQ-F-064.}

\requirement{REQ-F-055}{Опис огляду та лікування}{Лише призначений лікар повинен створювати опис огляду й лікування конкретного прийому. Причина звернення та результати огляду обов’язкові, діагноз і призначення — за потреби. Основні записи огляду й лікування пов’язані з прийомом.}{\hypertarget{AC-F-055}{\texttt{AC-F-055}}. Призначений лікар зберігає опис без діагнозу й призначень за наявності двох обов’язкових полів. Інший лікар і адміністратор не можуть створити цей опис. Збереження не завершує прийом автоматично.}

\requirement{REQ-F-056}{Результати аналізів та обстежень}{Будь-який ветеринар повинен мати можливість додавати результат аналізу чи обстеження окремо, без обов’язкового зв’язку з прийомом, у тому числі після його завершення. Обов’язкові назва й дата проведення, висновок необов’язковий; можна додати вкладення.}{\hypertarget{AC-F-056}{\texttt{AC-F-056}}. Інший ветеринар додає результат, який надійшов пізніше, без прийому та без висновку, але з назвою й датою. Відсутність назви чи дати відхиляється. Автоматичний обмін із лабораторією не потрібний.}

\requirement{REQ-F-057}{Формат і розмір вкладень}{Результат повинен підтримувати PDF, JPEG і PNG, до 5 файлів у кожній версії та до 10 000 000 байтів включно на файл. Перевіряється фактичний формат, не лише розширення.}{\hypertarget{AC-F-057}{\texttt{AC-F-057}}. П’ять дозволених файлів зберігаються, шостий відхиляється. Файл у 10 000 000 байтів дозволений, у 10 000 001 — ні. Файл іншого формату, перейменований на .pdf, відхиляється з поясненням.}

\requirement{REQ-F-058}{Узгоджене збереження результату та файлів}{Запис результату повинен зберігатися лише після успішного завантаження всіх обраних файлів. Лікар може повторити завантаження або прибрати проблемний файл. Невдале редагування не змінює збережену версію. Замінені чи вилучені з поточної версії файли залишаються у попередніх версіях.}{\hypertarget{AC-F-058}{\texttt{AC-F-058}}. При помилці другого файла нова версія не публікується. Після повторного успішного завантаження або явного вилучення проблемного файла версія зберігається. Попередню версію та її вкладення можна відкрити за належних прав.}

\requirement{REQ-F-059}{Редагування медичного запису}{Лише автор повинен редагувати збережений медичний запис, з обов’язковою причиною та збереженням попередньої версії. Для проведеного прийому редагування не може спорожнити причину звернення чи результати огляду. Блокування автора не передає права іншим; окремого коригувального механізму в цій версії немає.}{\hypertarget{AC-F-059}{\texttt{AC-F-059}}. Автор змінює запис із причиною; клієнт і ветеринари бачать попередню версію й причину. Інший автор або спроба залишити обов’язкове поле проведеного прийому порожнім отримує відмову.}

\requirement{REQ-F-060}{Помилково внесений медичний запис}{Замість остаточного видалення автор повинен мати можливість позначити запис помилковим із причиною. Він залишається в історії та не зараховується як коректний опис. Для проведеного прийому спочатку зберігається коректний опис; якщо прийому не було, спочатку виправляється його статус.}{\hypertarget{AC-F-060}{\texttt{AC-F-060}}. Без причини позначення не зберігається. Єдиний коректний опис проведеного прийому не можна зробити помилковим без заміни або попереднього виправлення статусу. Помилковий запис залишається видимим із позначкою.}

\requirement{REQ-F-061}{Виконані щеплення та обробки}{Виконане щеплення чи обробка повинні зберігатися медичним записом із видом процедури, датою виконання й назвою препарату. Дозу, серію та інші окремі додаткові поля для цього виду запису не включено. План наступної процедури є окремою рекомендацією.}{\hypertarget{AC-F-061}{\texttt{AC-F-061}}. Для процедури вашої клініки без будь-якого з трьох полів запис не зберігається. Запис виконаної процедури сам не обчислює наступну дату; рекомендацію визначає лікар.}

\requirement{REQ-F-062}{Історичні процедури з інших джерел}{Ветеринар повинен мати можливість додавати раніше виконані процедури, зокрема з іншої клініки, окремим медичним записом без прийому. Джерело обов’язкове; відсутні історичні дані явно позначаються як невідомі. Для процедур вашої клініки обов’язкові поля REQ-F-061 не послаблюються.}{\hypertarget{AC-F-062}{\texttt{AC-F-062}}. Історичне щеплення з джерелом «ветеринарний паспорт» та невідомими препаратом і точною датою зберігається. Без джерела — ні. Для поточної процедури клініки невідомий препарат не підміняє обов’язкову назву.}

\requirement{REQ-F-063}{Лікувальні призначення та дієта}{Лікар повинен зберігати лікувальні призначення й дієтичні рекомендації текстом у записі відповідного прийому. Окремого довідника препаратів та автоматичних розрахунків немає.}{\hypertarget{AC-F-063}{\texttt{AC-F-063}}. Призначений лікар зберігає текст лікування та дієти без вибору з довідника; клієнт бачить його після збереження, адміністратор — ні. Відсутність цих необов’язкових текстів не блокує профілактичний прийом.}

\subsection{Направлення та рекомендації}
\requirement{REQ-F-064}{Організаційна інформація адміністратора}{Адміністратор повинен бачити необхідність клінічного повторного прийому чи процедури, рекомендовану дату або строк та окреме поле «Рекомендація щодо запису». Це поле бачить також клієнт. Діагнози та лікування не включаються в організаційне подання. Домашні рекомендації не відкривають адміністратору медичний зміст.}{\hypertarget{AC-F-064}{\texttt{AC-F-064}}. Адміністратор бачить організаційну частину клінічної рекомендації та контакт клієнта, але не медичний текст чи вкладення. Автор заповнює окреме поле, а система не намагається автоматично виділити його з діагнозу.}

\requirement{REQ-F-065}{Направлення}{Ветеринар повинен мати можливість створити направлення з датою видачі, тим, що призначено, та напрямком звернення: аналіз, обстеження, спеціаліст або зовнішня установа. Конкретна установа й зв’язок із прийомом необов’язкові. Направлення не створює рекомендацію для запису автоматично.}{\hypertarget{AC-F-065}{\texttt{AC-F-065}}. Направлення до спеціаліста без назви конкретної установи й без прийому зберігається. Воно не з’являється у списку очікування. Для запису у клініці лікар окремо створює рекомендацію «У клініці».}

\requirement{REQ-F-066}{Авторство та формат рекомендації}{Будь-який ветеринар повинен мати можливість створити рекомендацію, обравши «У клініці» або «Вдома». Редагує лише автор; зберігаються попередня версія й обов’язкова причина виправлення. За пов’язаного підтвердженого або проведеного прийому зміна формату заборонена.}{\hypertarget{AC-F-066}{\texttt{AC-F-066}}. Інший ветеринар не редагує рекомендацію. Зміна дати автором із причиною зберігає версію. Переведення пов’язаної з підтвердженим прийомом рекомендації у домашню відхиляється.}

\requirement{REQ-F-067}{Дата рекомендації}{Лікар повинен указувати конкретну дату або строк у календарних днях від дати однієї явно обраної лікарем підстави: пов’язаного прийому або виконаної процедури. Відносний строк задається цілим числом календарних днів від 1 і більше; 0, від’ємні та дробові значення не допускаються. Верхня межа строку не встановлюється. Система обчислює дату й зберігає початковий строк. Подальша зміна дати базової події автоматично не змінює збережену дату рекомендації; за потреби автор виправляє рекомендацію з причиною та збереженням попередньої версії за REQ-F-066. Без такої підстави дозволена лише конкретна дата. Рекомендована дата є орієнтиром і не обмежує вибір іншого доступного дня. Автоматичних ветеринарних схем немає.}{\hypertarget{AC-F-067}{\texttt{AC-F-067}}. Строк 10 календарних днів від 28.09.2026 дає 08.10.2026 зі збереженням значення 10. Строк 1 день від 28.09.2026 дає 29.09.2026; значення 0, -1 і 1,5 відхиляються. Додатні цілі значення не відхиляються через верхню межу строку: її не встановлено. Без підстави відносний строк недоступний. Запис на інший вільний день дозволений. Якщо прийом був 01.10.2026, процедура — 03.10.2026, а строк — 10 днів, вибір процедури дає 13.10.2026, вибір прийому — 11.10.2026. Виправлення дати процедури на 04.10.2026 не змінює збережену рекомендацію на 13.10.2026 без явного виправлення автором.}

\requirement{REQ-F-068}{Список очікування запису}{Система повинна показувати адміністратору чинні рекомендації «У клініці», для яких немає підтвердженого або проведеного прийому та дозволено запис тварини. Рядок містить клієнта, контакт, тварину, дату або строк, прийом чи процедуру та рекомендацію щодо запису. Домашні й скасовані рекомендації не включаються.}{\hypertarget{AC-F-068}{\texttt{AC-F-068}}. Клінічна рекомендація без прийому видима з усіма переліченими полями, домашня — ні. Минулість рекомендованої дати сама не вилучає її. Смерть, припинення опіки й скасування рекомендації вилучають її з цього подання.}

\requirement{REQ-F-069}{Запис за рекомендацією}{Адміністратор і клієнт повинні мати можливість почати запис із конкретної рекомендації. Лише успішне збереження встановлює зв’язок і прибирає її зі списку. За однією рекомендацією не може бути більш ніж один прийом у сукупності статусів «Підтверджено» та «Проведений», включно з конкурентними запитами.}{\hypertarget{AC-F-069}{\texttt{AC-F-069}}. Відкриття форми та невдале бронювання не прибирають рекомендацію. Успішний запис прибирає її; повторний конкурентний запис за тією самою рекомендацією на інший час не створюється. За проведеного прийому повторний запис за нею недоступний.}

\requirement{REQ-F-070}{Зв’язок рекомендації зі статусом прийому}{Скасування або неявка повинні повертати чинну клінічну рекомендацію до списку, навіть із минулою датою, якщо стан тварини дозволяє запис. Перенесення зберігає зв’язок. Проведення не встановлює окремого стану рекомендації «Виконана». Виправлення «Проведений» до «Підтверджено» не повертає її до списку.}{\hypertarget{AC-F-070}{\texttt{AC-F-070}}. Скасований прийом повертає допустиму рекомендацію; проведений і виправлений до підтвердженого — ні. За смерті, припинення опіки або скасованої рекомендації повернення не відбувається.}

\requirement{REQ-F-071}{Скасування рекомендації}{Автор повинен мати можливість скасувати рекомендацію з обов’язковою причиною. Вона зберігається в історії, не очікує запису й не дозволяє нового бронювання за нею. Зміна або скасування рекомендації не змінює час, лікаря чи статус уже створеного прийому. Цей прийом можна переносити за REQ-F-043 і REQ-F-044 зі збереженням зв’язку та скасованого стану рекомендації.}{\hypertarget{AC-F-071}{\texttt{AC-F-071}}. Після скасування рекомендації чинний прийом залишається без змін; його подальше скасування або неявка не повертають рекомендацію до списку. Без причини скасування рекомендації відхиляється.}

\requirement{REQ-F-072}{Домашні процедури}{Для домашньої процедури система повинна зберігати інструкцію та рекомендовану дату. Онлайн-запис за нею та позначка виконання клієнтом не передбачені.}{\hypertarget{AC-F-072}{\texttt{AC-F-072}}. Клієнт переглядає домашню рекомендацію; у ній немає дії запису або відмітки виконання, а в адміністративному списку її немає.}

\subsection{Поведінка зовнішніх сервісів}
\requirement{REQ-F-073}{Помилка SMS-сервісу}{При відхиленні надсилання або відсутності відповіді SMS-сервісу протягом 10 секунд система повинна повідомляти, що надсилання не вдалося або не підтверджене. Без перевірки коду реєстрація, зміна номера й відновлення не завершуються. Повторний запит підлягає правилу 60 секунд. Прийняття запиту сервісом не означає доставки телефону.}{\hypertarget{AC-F-073}{\texttt{AC-F-073}}. Імітація помилки та тайм-ауту не підтверджує номер і не змінює доступ; користувач отримує відповідний стан і можливість повтору за строком. Позитивна відповідь сервісу сама не завершує дію.}

\requirement{REQ-F-074}{Сповіщення про невдале резервне копіювання}{Система повинна надсилати електронний лист відповідальній технічній особі на адресу в налаштуваннях розгортання протягом 5 хвилин після виявлення помилки резервного копіювання. Нова роль у застосунку для цієї особи не створюється.}{\hypertarget{AC-F-074}{\texttt{AC-F-074}}. Штучно спричинена помилка створює лист на налаштовану адресу не пізніше ніж через 5 хвилин від виявлення. Надсилання фіксується у перевірці; збій самого поштового сервісу залишається OPEN-006.}

\section{Нефункціональні вимоги}
\subsection{Інтерфейс і часові правила}
\requirement{REQ-NF-001}{Український адаптивний інтерфейс}{Інтерфейс повинен бути лише українською та адаптуватися до ПК і смартфона. Працівники працюють через браузер на ПК, клієнти — на ПК або смартфоні. Окремих мобільних застосунків немає.}{\hypertarget{AC-NF-001}{\texttt{AC-NF-001}}. На погоджених розмірах екрана з DEP-001 основні потоки доступні без обрізаних полів і недоступних дій; повідомлення та підписи українською. Остаточні розміри екрана фіксуються до приймальної перевірки.}

\requirement{REQ-NF-002}{Підтримувані браузери}{Система повинна працювати у Chrome, Edge та Firefox на ПК, Chrome на Android і Safari на iOS. Конкретні версії фіксуються під час прийняття.}{\hypertarget{AC-NF-002}{\texttt{AC-NF-002}}. Протокол приймання містить точні ОС та версії всіх п’яти комбінацій браузера/платформи й результати основних потоків; неперевірені версії не оголошуються підтриманими без доказу.}

\requirement{REQ-NF-003}{Часовий пояс і межі інтервалів}{Розклад, час прийомів і рекомендовані дати повинні показуватися за Europe/Kyiv незалежно від пристрою. Інтервали прийомів трактуються як [початок, кінець), тому спільна межа не є перетином. Часові умови перевіряються системою, а не лише інтерфейсом. Неіснуючі або неоднозначні значення місцевого часу під час сезонного переведення годинника не приймаються й не пропонуються; система показує пояснення без автоматичного зсуву часу.}{\hypertarget{AC-NF-003}{\texttt{AC-NF-003}}. Зміна часового поясу пристрою не змінює показаного часу клініки. Прийоми 10:00–10:20 та 10:20–10:40 дозволені. Неіснуюче або неоднозначне значення місцевого часу не пропонується для вибору, а його пряме введення відхиляється з поясненням; автоматичної підстановки іншого часу немає.}

\subsection{Продуктивність та інтеграції}
\requirement{REQ-NF-004}{Швидкодія основних операцій}{За 50 одночасно активних користувачів пошук, відкриття картки та збереження запису на прийом повинні завершуватися не довше ніж за 2 секунди у щонайменше 95 відсотках операцій кожного типу окремо. Тестові дані: 5000 клієнтів, 10 000 тварин, 100 000 прийомів, 200 000 медичних записів. До 5 адміністраторів і 10 лікарів є орієнтиром, не лімітом реєстрації.}{\hypertarget{AC-NF-004}{\texttt{AC-NF-004}}. У погодженому середовищі DEP-001 окремі вимірювання кожного типу дають частку операцій до 2 секунд включно не менше 95 відсотків. Протокол містить обсяги, навантаження, тривалість прогону, кількість операцій і помилки; конкретний профіль прогону погоджується до тесту.}

\requirement{REQ-NF-005}{Швидкість завантаження файла}{Один дозволений файл у 10 000 000 байтів повинен завантажуватися не довше ніж за 15 секунд у 95 відсотках спроб за стабільної вихідної швидкості клієнта не менше 10 Мбіт/с та до 50 активних користувачів у визначеному середовищі.}{\hypertarget{AC-NF-005}{\texttt{AC-NF-005}}. Протокол DEP-001 фіксує швидкість мережі, кількість спроб та час від початку завантаження до його підтвердження застосунком; щонайменше 95 відсотків спроб вкладаються у 15 секунд. Помилки подаються окремо, а не маскуються як швидкі успіхи.}

\requirement{REQ-NF-006}{Граничний час відповіді SMS-сервісу}{Очікування відповіді SMS-сервісу про прийняття чи відхилення запиту не повинно перевищувати 10 секунд. Доставка на телефон є окремою інтеграційною перевіркою, без погодженого числового строку доставки.}{\hypertarget{AC-NF-006}{\texttt{AC-NF-006}}. За відсутності відповіді через 10 секунд застосунок припиняє очікування та відображає непідтверджене надсилання за REQ-F-073. Окремий тест із реальним сервісом засвідчує отримання коду на контрольному українському номері.}

\subsection{Резервування і відновлення}
\requirement{REQ-NF-007}{Резервування та допустима втрата}{Дані й вкладення повинні копіюватися автоматично щонайменше щогодини. Допустима втрата після серйозного збою — не більше 1 години; копія повинна бути придатною до узгодженого відновлення записів і файлів.}{\hypertarget{AC-NF-007}{\texttt{AC-NF-007}}. За результатами резервування та пробного відновлення втрата змін не перевищує години від моменту збою. Посилання на вкладення відновлених записів відкривають відповідні файли. Сам запуск невдалого копіювання не підтверджує виконання вимоги.}

\requirement{REQ-NF-008}{Зберігання резервних копій}{Резервні копії повинні зберігатися 7 діб окремо від основного сервера застосунку. Постачальник і технологія ще не визначені.}{\hypertarget{AC-NF-008}{\texttt{AC-NF-008}}. Наявні придатні копії за погоджений семиденний період; відновлення виконується без доступу до основного сервера. Конфігурація розгортання підтверджує окреме розміщення. Це не строк зберігання медичної історії.}

\requirement{REQ-NF-009}{Відновлення за чотири години}{Відповідальна за технічне обслуговування особа повинна відновити роботу з резервної копії не пізніше 4 годин від припинення роботи через серйозний збій. Короткі мережеві помилки не є цим сценарієм.}{\hypertarget{AC-NF-009}{\texttt{AC-NF-009}}. Перед прийняттям пробне відновлення в окремому середовищі завершується в межах 4 годин від модельованої зупинки, забезпечує доступ до карток і файлів та виконує REQ-NF-007. Протокол часу й перевірок збережено.}

\subsection{Безпека}
\requirement{REQ-NF-010}{Правила пароля}{Паролі клієнтів і працівників повинні мати 15–128 кодових точок Unicode; пробіли дозволені й не обрізаються. Введений пароль автоматично не змінюється. Обов’язкового поєднання великих літер, цифр і спеціальних символів немає.}{\hypertarget{AC-NF-010}{\texttt{AC-NF-010}}. Паролі довжиною 15 і 128, у тому числі з пробілами, приймаються; 14 і 129 — відхиляються. Пароль із дозволених символів не відхиляється лише через відсутність цифр або великих літер. Літера та окремий комбінований наголос рахуються як дві кодові точки, незалежно від відображення одним знаком. Початкові й кінцеві пробіли зберігаються; пароль не нормалізується автоматично.}

\requirement{REQ-NF-011}{Захист входу від повторних помилок}{Після 5 послідовних неправильних паролів для облікового запису вхід повинен блокуватися на 15 хвилин. Спроби під час блокування не подовжують строк. Лічильник обнуляється після успішного входу або завершення блокування. Повідомлення не розкриває існування логіна чи телефону.}{\hypertarget{AC-NF-011}{\texttt{AC-NF-011}}. Після п’ятої помилки навіть правильний пароль до завершення 15 хвилин не надає входу; спроба всередині строку його не подовжує. Після завершення правильний пароль працює. Для неіснуючих та неправильних даних повідомлення не розрізняє причину.}

\requirement{REQ-NF-012}{Завершення сеансів через бездіяльність}{Сеанс працівника повинен завершуватися після 30 хвилин бездіяльності, клієнта — після 24 годин. Активність — явна дія користувача з запитом до системи; фонова робота, відкрита вкладка й введення тексту без надсилання не продовжують сеанс. Потрібний повторний вхід, незбережені зміни не зберігаються автоматично.}{\hypertarget{AC-NF-012}{\texttt{AC-NF-012}}. Після відповідного строку без запиту користувача захищена дія потребує входу. Фонові запити та введення довгого опису не змінюють результат; пошук або відкриття картки продовжує строк. Автозбереження незбереженого тексту не відбувається.}

\requirement{REQ-NF-013}{Захищене передавання}{Увесь застосунок повинен використовувати HTTPS для входу, карток, вкладень та інших операцій. Незахищене HTTP не передає ці дані.}{\hypertarget{AC-NF-013}{\texttt{AC-NF-013}}. Мережеві перевірки основних потоків і прямих запитів не виявляють передавання паролів, карток або файлів через HTTP. HTTPS-конфігурацію перевірено у цільовому середовищі.}

\requirement{REQ-NF-014}{Зберігання секретів}{Паролі повинні зберігатися лише як стійкі односторонні хеші, не у відкритому або відновлюваному вигляді. Паролі, SMS-коди й секретні частини одноразових посилань не повинні потрапляти до журналів. Алгоритм і параметри обираються при проєктуванні реалізації.}{\hypertarget{AC-NF-014}{\texttt{AC-NF-014}}. Огляд сховища та реалізації не виявляє відкритих або розшифровуваних паролів. Після тестів із контрольними секретами їх немає у журналах.}

\requirement{REQ-NF-015}{Серверне розмежування доступу}{Система повинна перевіряти права на сервері для кожного запиту, включно з файлами, попередніми версіями та підстановкою ідентифікатора. Прихована кнопка не є достатнім захистом. Пряме посилання не надає додаткових прав.}{\hypertarget{AC-NF-015}{\texttt{AC-NF-015}}. Сторонній клієнт не читає й не змінює чужу картку; адміністратор не читає медичний текст або файл; заблокований працівник не виконує захищених дій. Негативні перевірки прямими запитами проходять для актуальних і попередніх версій.}

\section{Вимоги до процесу розробки}
Цей розділ визначає роботу розробника й coding agent, а не бізнес-функції клініки. Наявність історії медичних виправлень не виконує вимогу журналу TDD, і навпаки. Методи та матеріали перевірки цього розділу — репозиторій, журнали, дифи й рішення людини, а не клієнтський інтерфейс.

\subsection{Development Process Requirements}
\requirement{REQ-DP-001}{Погоджена вимога перед реалізацією}{Зміна функціональної поведінки повинна починатися з погодженої вимоги зі стабільним ID та критерієм прийняття. Coding agent не додає суттєву поведінку для закриття невизначеності самостійно.}{\hypertarget{AC-DP-001}{\texttt{AC-DP-001}}. Для перевірюваної зміни до початку реалізації існують погоджені REQ і AC; зміни обсягу мають окреме людське рішення. Чернетка до SPEC-GATE не оголошується затвердженою базовою версією.}

\requirement{REQ-DP-002}{Обов’язковий TDD}{Кожна зміна функціональної поведінки повинна проходити цикл: автоматизований тест до реалізації; запуск і падіння з очікуваної причини; реалізація до проходження; за потреби рефакторинг лише після успіху та повторний запуск. Суто текстові зміни документації не потребують цього циклу.}{\hypertarget{AC-DP-002}{\texttt{AC-DP-002}}. Для вибраної зміни є впорядковані докази Red і Green, а після рефакторингу — повторного Green. Падіння через помилку середовища не підміняє очікуване падіння перевірки поведінки.}

\requirement{REQ-DP-003}{Логи доказів тестування}{У /logs повинні зберігатися ID вимоги, команди запуску тестів, результат до реалізації з причиною очікуваного падіння та результат після реалізації; якщо був рефакторинг — результат після нього.}{\hypertarget{AC-DP-003}{\texttt{AC-DP-003}}. За ID вимоги можна знайти всі потрібні запуски, команди та результати. Відсутність доказу до реалізації не маскується пізнішим успішним запуском.}

\requirement{REQ-DP-004}{Журнал роботи coding agent}{У /logs повинен зберігатися стислий запис завдання, пов’язаних ID вимог, запропонованих агентом змін, виконаних перевірок і рішення людини: прийнято, виправлено або відхилено, з поясненням. Повні діалоги не обов’язкові.}{\hypertarget{AC-DP-004}{\texttt{AC-DP-004}}. Для зміни агента журнал дозволяє встановити завдання, REQ, внесок, перевірки та людське рішення. Він відокремлений від історії змін медичних даних і не містить секретів.}

\requirement{REQ-DP-005}{Версіювання та простежуваність}{Базова специфікація повинна зберігатися як LaTeX у /spec. ID не перенумеровуються через перестановку; вилучені ID не використовуються повторно. Для цієї лабораторної роботи матриця пов’язує кожну вимогу з критерієм прийняття та BDD-сценарієм, якщо він є. Обов’язкове майбутнє заповнення посилань на тікети, код і документацію не встановлюється. Зв’язок тестів із вимогами забезпечується в межах TDD.}{\hypertarget{AC-DP-005}{\texttt{AC-DP-005}}. Кожен REQ має унікальний ID, AC і рядок матриці. Усі BDD-посилання існують. Історія Git дозволяє порівняти версії; відсутність BDD не приховує відсутності способу перевірки.}

\requirement{REQ-DP-006}{Підсумок SKED та окремий SPEC-GATE}{Перед затвердженням повинні бути збережені підсумок SKED у /logs із невизначеностями, людськими рішеннями й виправленими припущеннями, а також виконана окрема контрольна перевірка специфікації. Результат SPEC-GATE — PASS або FAIL — фіксується після контрольної перевірки. Після FAIL опрацьовуються блокуючі зауваження та проводиться повторна перевірка. PASS дозволяє подати специфікацію на людське затвердження, але не затверджує її автоматично. Коміт затвердженої baseline створюється лише після людського затвердження.}{\hypertarget{AC-DP-006}{\texttt{AC-DP-006}}. Перевірка окремо оцінює суперечності, терміни, перевірюваність і простежуваність; рішення має фактичний статус. Після перевірки збережено її фактичний результат PASS або FAIL. За FAIL передбачено опрацювання блокуючих зауважень і повторну перевірку; за PASS без людського затвердження немає затвердженої baseline або коміту її затвердження. Коміт затвердженої baseline має попереднє людське затвердження після PASS.}

\section{BDD сценарії Given–When–Then}
Наведені приклади покривають ключові потоки та суттєві негативні й граничні випадки. Дані в сценаріях є синтетичними. Зазначені окремі прогони є варіантами одного правила; при автоматизації вони стають окремими тестовими випадками. Перелік не замінює всіх AC і не заявляє, що тести вже реалізовані.

\scenario{BDD-001}{Реєстрація без тварини}{Новий клієнт має вільний український номер, заповнені ім’я та прізвище й допустимий пароль; тварини та пошти немає.}{Клієнт підтверджує номер правильним чинним SMS-кодом і завершує реєстрацію.}{Створено один профіль без тварин; доступ дозволений, додавання тварини не є умовою реєстрації.}

\scenario{BDD-002}{Існуючий номер не створює дубля}{Адміністратор уже створив профіль із номером +380671234567.}{Інша реєстрація використовує +380 (67) 123-45-67.}{Новий профіль не створюється, а наявний не стає доступним лише через збіг номера.}

\scenario{BDD-003}{Межі й заміна коду}{Код A виданий для реєстрації о 10:00:00, код ще не використаний.}{Користувач запитує повтор о 10:00:59, потім о 10:01:00; система видає код B, після чого користувач вводить A.}{Перший повтор заборонений, другий дозволений; A відхиляється. B не підходить для зміни номера й після завершення власних 5 хвилин не приймається. Окремий прогін п’яти неправильних введень вичерпує спроби.}

\scenario{BDD-004}{Одноразове запрошення клієнта}{Профіль без доступу створений адміністратором, посилання чинне менше 24 годин.}{Клієнт установлює пароль за посиланням і повторно відкриває його для встановлення іншого пароля.}{Перша операція надає доступ до того самого профілю, друга відхиляється. В окремих прогонах прострочення і перевидача посилання не видаляють профіль, але роблять старе посилання непридатним.}

\scenario{BDD-005}{Тимчасове блокування входу}{Клієнт має чинний пароль і нуль невдалих спроб.}{П’ять разів поспіль введено неправильний пароль, потім правильний під час 15-хвилинного блокування.}{Вхід не надається; спроба не подовжує строк. Після його завершення правильний пароль надає вхід. Текст помилки не розкриває існування номера; успішний вхід обнуляє лічильник.}

\scenario{BDD-006}{Відновлення пароля клієнта}{У клієнта є два активні сеанси й доступ до зареєстрованого номера.}{Клієнт підтверджує код відновлення й установлює допустимий новий пароль.}{Обидва попередні сеанси завершені, старий пароль не працює, потрібен новий вхід. Повторне використання коду або код реєстрації не дозволяє ще одну зміну.}

\scenario{BDD-007}{Зміна номера без часткового результату}{Клієнт увійшов зі старим номером; новий номер вільний.}{Він уводить поточний пароль, але ще не підтверджує новий номер, а потім успішно підтверджує SMS-код.}{До підтвердження чинний старий номер; після успіху — новий, усі старі сеанси завершені. В окремому прогоні зайнятий новий номер залишає старі дані без змін.}

\scenario{BDD-008}{Особисте відновлення без копії документа}{Клієнт утратив кабінет і старий номер; адміністратор перевірив документ наживо.}{Адміністратор фіксує позитивний результат, клієнт підтверджує новий номер і встановлює пароль за посиланням.}{Лише після двох кроків доступ змінено й попередні сеанси завершено. Збережені адміністратор, час і результат без копії й реквізитів документа. Негативний результат в окремому прогоні не запускає зміну. В окремих прогонах посилання діє 24 години від створення; успішно підтверджений номер залишається підтвердженим до цієї межі без окремого таймера. Повторне використання, прострочене посилання або підтвердження іншого профілю, номера чи операції не дають змінити доступ. Нове посилання скасовує попередню операцію й потребує нового підтвердження номера; повтор SMS замінює лише код. Якщо номер зайнято до спільного збереження, старі дані доступу залишаються незмінними.}

\scenario{BDD-009}{Робочий доступ окремий від особистого}{Перший адміністратор створений при налаштуванні; ветеринар також має клієнтський профіль.}{Адміністратор створює робочий логін vet.ivan і передає одноразове посилання; працівник установлює пароль.}{Робочий вхід окремий від клієнтського; роль одна й незмінна. Логін Vet.Ivan повторно не створюється. Відновлення через нове посилання не передає пароль адміністратору й завершує попередні сеанси.}

\scenario{BDD-010}{Блокування ветеринара з майбутніми прийомами}{У ветеринара є активний сеанс, збережені медичні записи й майбутній прийом.}{Адміністратор блокує його, а згодом розблоковує.}{Під час блокування сеанс недійсний, нові майбутні записи до лікаря недоступні, наявний прийом і авторство збережені. Після розблокування потрібний новий вхід.}

\scenario{BDD-011}{Останній активний адміністратор}{Адміністратор A є єдиним незаблокованим адміністратором з установленим паролем; B ще не встановив пароль.}{Надходить запит заблокувати A або A намагається заблокувати себе.}{Дія відхиляється; B не зараховується як працездатна заміна. Перевірка паралельних блокувань також не залишає систему без активного адміністратора.}

\scenario{BDD-012}{Непідтверджена тварина та перший запис}{Клієнт підтвердив наявність паспорта, але номер невідомий; задані кличка, вид і стать «Невідомо», народження лише роком.}{Він додає тварину та бронює її перший доступний прийом.}{Картка створена як непідтверджена, день і місяць народження не вигадані; запис підтверджений. До першої перевірки клієнт редагує основні дані. В окремих прогонах майбутня точна дата народження та рік, пізніший за поточний рік Europe/Kyiv, відхиляються; поточні дата або рік допустимі.}

\scenario{BDD-013}{Конкурентний дубль мікрочипа}{Дві різні картки ще не мають номера 001234567890123.}{Два запити одночасно намагаються зберегти цей номер, у тому числі через схвалення пропозиції змін.}{Лише один зберігає номер із початковими нулями. Інший отримує відмову та пораду звернутися до адміністратора без відомостей іншого власника.}

\scenario{BDD-014}{Первинне підтвердження не блокує допомогу}{Клієнтська картка непідтверджена й даних поки недостатньо.}{Працівник залишає пояснення, а призначений лікар зберігає опис і позначає прийом проведеним після його початку.}{Картка лишається непідтвердженою, але прийом проведений. Пізніше працівник може виправити дані й підтвердити їх. Окремо створена працівником при особистому зверненні картка без паспорта одразу підтверджена.}

\scenario{BDD-015}{Пропозиція не підміняє чинних даних}{Картка підтверджена, чинна кличка «Мурка».}{Клієнт пропонує іншу кличку, пробує подати другу пропозицію, а працівник відхиляє першу з причиною.}{До рішення в картці чинна «Мурка», друга пропозиція заблокована; після відхилення причина видима й можна подати нову. Альтернативне схвалення всього набору полів замінює чинні значення.}

\scenario{BDD-016}{Виправлення працівника робить пропозицію неактуальною}{Нерозглянута пропозиція змінює кличку й породу.}{Ветеринар безпосередньо виправляє породу.}{Усю пропозицію відхилено як неактуальну з причиною для клієнта. В окремому прогоні зміна лише номера паспорта не закриває пропозицію.}

\scenario{BDD-017}{Перше рішення щодо пропозиції}{Адміністратор і ветеринар відкрили ту саму нерозглянуту пропозицію.}{Вони одночасно надсилають схвалення й відхилення.}{Зберігається лише перше успішне рішення; другий працівник отримує повідомлення й актуальний стан без перезапису.}

\scenario{BDD-018}{Наслідки смерті або припинення опіки}{Тварина має власника, відкриту пропозицію змін, чинну рекомендацію та прийоми на 09:00, 10:00 і 11:00.}{Працівник о 10:00 вносить одну з подій, що сталася о 09:00.}{Запис на 11:00 скасований з причиною, два інші не змінені автоматично. Пропозиція відхилена, рекомендація не в списку, нові клієнтські записи й зміни заборонені. Картка в окремому розділі й профіль клієнта зберігаються. Сценарій виконується окремо для обох подій. В окремому прогоні з непідтвердженою карткою пряме редагування основних даних клієнтом також заборонене; перегляд зберігається. В окремих прогонах внесення події з майбутніми датою/часом або майбутньою датою без часу відхиляється.}

\scenario{BDD-019}{Гонка між записом і позначкою події}{Клієнт бронює майбутній прийом за рекомендацією одночасно з фіксацією смерті працівником.}{Система зберігає спочатку запис або спочатку смерть у двох окремих прогонах.}{У першому прогоні запис потім скасований подією, у другому — створення відхилене. В обох рекомендація не очікує запису й немає чинного майбутнього прийому.}

\scenario{BDD-020}{Виправлення незалежних позначок}{Зафіксовано і смерть, і припинення опіки; майбутні записи раніше скасовано.}{Працівник із причиною виправляє смерть, потім окремо помилкове припинення опіки.}{Після першого виправлення доступ до нового запису не з’являється. Після відновлення поточного власника допустимі чинні рекомендації без підтвердженого або проведеного прийому повертаються; скасовані прийоми не відновлюються.}

\scenario{BDD-021}{Пошук з архівними станами}{Є однойменні клієнти з різними телефонами та тварини «Мурка», включно з припиненою опікою.}{Працівник шукає частину імені, «мур» і точний номер мікрочипа.}{Частковий текст ігнорує регістр; клієнти розрізняються телефонами. Архівна картка має позначку й текст «Поточний власник відсутній». Точний мікрочип знаходить одну картку; частина номера не є точним збігом.}

\scenario{BDD-022}{Формування двадцятихвилинних інтервалів}{Адміністратор задав на дату робочий інтервал 09:00–09:50 за Europe/Kyiv.}{Клієнт відкриває доступний час із пристрою з іншим часовим поясом.}{Показано 09:00 та 09:20 за часом клініки, але не 09:40. Суміжні двадцятихвилинні прийоми не конфліктують. В окремих прогонах неіснуючий та неоднозначний місцевий час сезонного переходу не пропонується, а пряме введення відхиляється з поясненням без автоматичного зсуву.}

\scenario{BDD-023}{Конфлікт зміни розкладу}{Є підтверджений майбутній прийом 10:00–10:20.}{Адміністратор змінює робочий інтервал так, що прийом більше не вміщується.}{Зміна не зберігається, показано запис-перешкоду. Після його окремого перенесення або скасування зміна може бути збережена. В окремому прогоні зміна початку робочого інтервалу з 09:00 на 09:10 за незмінного кінця 13:00 дозволена: прийом 10:00–10:20 залишається, а нові інтервали, що перетинають його, недоступні. Історичний прийом поза розкладом не є перешкодою.}

\scenario{BDD-024}{Негайне підтвердження онлайн-запису}{Клієнт має доступну для запису тварину; вільний прийом лікаря починається о 10:20, зараз 10:00.}{Клієнт обирає лікаря й 10:20 та зберігає запис.}{Статус одразу «Підтверджено», без погодження адміністратора. В окремому прогоні створення клієнтом запису на 10:00 відхиляється.}

\scenario{BDD-025}{Один час для двох користувачів}{Два користувачі бачать вільний час одного лікаря 10:00–10:20.}{Вони одночасно зберігають різні записи.}{Лише перший успішно збережений запис підтверджено, іншому показано зайнятість; подвійного бронювання немає.}

\scenario{BDD-026}{Перетин прийомів однієї тварини}{Тварина має прийом 10:00–10:20 до лікаря A; перевірка виконується окремо для статусів «Підтверджено» та «Проведений».}{Адміністратор пробує додати їй історичний прийом 10:10–10:30 до лікаря B.}{Створення відхиляється незалежно від вільності B та екстреної причини; суміжний інтервал 10:20–10:40 не порушує саме це правило.}

\scenario{BDD-027}{Повтор після втраченої відповіді}{Запис успішно збережений, але відповідь користувачу втрачена через обрив зв’язку.}{Надходить повтор тієї самої спроби створення.}{Повертається підтвердження того самого запису; немає дубля чи повідомлення про зайнятість власним записом.}

\scenario{BDD-028}{Невдале й успішне перенесення}{Клієнт має майбутній запис за рекомендацією, а новий обраний час уже зайнято.}{Він пробує перенести запис, а потім обирає інший вільний час і повторює.}{Перша спроба залишає старий запис незмінним; друга змінює час і лише тоді звільняє старий. Зв’язок із рекомендацією збережений. Після початку первісного прийому перенесення заборонене також адміністратору. Перенесення на минулий час відхиляється і для клієнта, і для адміністратора. В окремому прогоні скасована рекомендація не заважає успішному перенесенню на майбутній час; її зв’язок і скасований стан зберігаються.}

\scenario{BDD-029}{Зміна лікаря за наявності медичного запису}{До майбутнього прийому вже прив’язаний медичний запис призначеного лікаря.}{Клієнт для своєї тварини або адміністратор у двох окремих прогонах намагається перенести його до іншого ветеринара.}{Зміна лікаря заблокована, авторство й запис незмінні. В окремих прогонах без медичного запису обом ролям дозволене перенесення до доступного лікаря на майбутній вільний час; якщо новий час недоступний, попередній запис залишається без змін.}

\scenario{BDD-030}{Скасування клієнтом і клінікою}{Є підтверджений прийом, пов’язаний із чинною клінічною рекомендацією.}{Клієнт скасовує його до початку; в окремому прогоні після початку це робить адміністратор через проблему клініки з причиною.}{Скасування успішне й причина клініки видима клієнту; рекомендація повертається за виконаних умов. Клієнт після початку та адміністратор без причини після початку не можуть скасувати прийом.}

\scenario{BDD-031}{Окрема дія проведення прийому}{Час прийому настав; призначений лікар ще не зберіг обов’язкового опису.}{Він пробує встановити «Проведений», потім зберігає причину звернення й результати огляду та повторює дію.}{Перша спроба відхилена; саме збереження опису залишає «Підтверджено»; лише повторна окрема дія дає «Проведений». Інший лікар та адміністратор цю дію не виконують.}

\scenario{BDD-032}{Неявка повертає рекомендацію}{Час початку підтвердженого прийому минув; рекомендація чинна, тварина доступна для запису.}{Адміністратор установлює «Неявка».}{Рекомендація знову в списку навіть із минулою рекомендованою датою. Окремий прогін до початку відхиляє позначення неявки. Після позначення неявки спроба створити інший історичний прийом до того самого лікаря з перетином цього інтервалу відхиляється.}

\scenario{BDD-033}{Відновлення майбутнього скасованого запису}{Запис помилково скасований, але первісний час уже зайнятий іншим прийомом.}{Адміністратор пробує відновити його з причиною.}{Відновлення заблоковане, інший прийом не скасовується. За вільного часу й чинного розкладу окремий прогін повертає «Підтверджено» та зберігає автора, час, причину й обидва статуси. В окремих прогонах навіть за вільного часу та чинного розкладу відновлення майбутнього запису відхиляється для померлої тварини, за відсутності поточного власника або до заблокованого лікаря.}

\scenario{BDD-034}{Виправлення історичного статусу}{Минулий прийом помилково позначений неявкою, конфліктів немає, відповідного інтервалу в актуальному розкладі вже немає.}{Адміністратор виправляє неявку з причиною.}{Запис повертається до «Підтверджено»; тільки призначений лікар окремо визначає проведення. Якщо вже існує інший підтверджений або проведений прийом за рекомендацією, відновлення блокується. Виправлення лікарем «Проведений» не повертає рекомендацію в список. В окремому прогоні заблокований лікар не перешкоджає відновленню минулого запису за допустимої хронології REQ-F-052; прийом, що не відповідає цій хронології, не відновлюється.}

\scenario{BDD-035}{Минулий прийом поза розкладом до заблокованого лікаря}{Прийом фактично відбувся раніше, не перетинає інших, але час був поза розкладом і лікар тепер заблокований.}{Адміністратор вносить фактичний минулий початок.}{Створено двадцятихвилинний «Підтверджений» прийом без обмеження давності. Інший лікар не може додати його опис чи завершити; потрібне можливе розблокування призначеного лікаря.}

\scenario{BDD-036}{Історичний прийом відносно події тварини}{Опіку припинено о 10:20; в окремому прогоні смерть настала о 10:10.}{Адміністратор вносить прийом 10:00–10:20.}{Обидва варіанти дозволені. Для припинення опіки прийом 10:10–10:30 відхиляється. За невідомого часу події того самого дня потрібні явне підтвердження допустимості, адміністратор, час і причина.}

\scenario{BDD-037}{Розділення медичного й організаційного доступу}{Лікар зберіг діагноз, дієту й клінічну рекомендацію з окремою вказівкою щодо запису.}{Картку відкривають належний клієнт, інший ветеринар і адміністратор.}{Клієнт і ветеринар бачать медичний запис; адміністратор бачить організаційну частину рекомендації, але не діагноз, дієту та медичні версії.}

\scenario{BDD-038}{Пряма адреса файла не обходить права}{Є медичне вкладення й попередня версія, доступні власнику картки та ветеринарам.}{Сторонній клієнт і адміністратор надсилають прямі запити з ідентифікаторами цих ресурсів.}{Обидва отримують відмову без вмісту файлів і версій; приховування кнопки не є єдиною перевіркою. Належний клієнт і ветеринар мають доступ.}

\scenario{BDD-039}{Пізній результат без прийому}{Огляд завершений; результат зовнішнього аналізу отримано пізніше іншим ветеринаром.}{Ветеринар додає назву, дату проведення й файл без прив’язки до прийому та без висновку.}{Результат збережено й видно клієнту; інтеграція з лабораторією не потрібна. Без назви або дати звичайного результату збереження відхиляється.}

\scenario{BDD-040}{Межі вкладень і часткова помилка}{Лікар додає дозволені PDF, JPEG або PNG, кожен не більше 10 000 000 байтів.}{Один файл не завантажується; окремо перевіряються шостий файл, файл 10 000 001 байт і підроблене розширення.}{Новий результат не зберігається до успіху всіх обраних файлів або явного вилучення проблемного; три порушення відхиляються. П’ять дозволених файлів, включно з граничним розміром, зберігаються.}

\scenario{BDD-041}{Редагування з історією вкладень}{Результат автора A має текст і файл у збереженій версії.}{A змінює текст і замінює файл із причиною.}{Поточна версія оновлена, попередній текст і файл доступні належному клієнту та ветеринарам. Редагування іншою особою або без причини відхиляється; невдале завантаження не змінює чинної версії.}

\scenario{BDD-042}{Помилковий опис проведеного прийому}{Проведений прийом має один коректний опис.}{Автор пробує зробити його помилковим без заміни, потім зберігає коректний новий опис і повторює з причиною.}{Перша дія відхилена, друга залишає старий опис у історії як помилковий. Якщо прийому взагалі не було, перед позначкою виправляється статус. Очищення обов’язкових полів проведеного прийому редагуванням також заборонене.}

\scenario{BDD-043}{Виконана та історична процедура}{Лікар документує процедуру клініки та окремо щеплення з давнього паспорта.}{Для першої вводить вид, дату й препарат, для другої — джерело з явно невідомим препаратом.}{Обидва записи зберігаються. Без препарату перша не зберігається, без джерела друга не зберігається. Наступна дата не обчислюється автоматично за схемою.}

\scenario{BDD-044}{Зовнішнє направлення без бронювання}{Ветеринар направляє тварину на обстеження поза клінікою.}{Він указує дату видачі, обстеження й напрямок без конкретної установи та без прийому.}{Направлення збережено, автоматичного обміну та рекомендації для запису немає; до списку адміністратора воно не потрапляє.}

\scenario{BDD-045}{Зміна й скасування рекомендації не змінюють бронювання}{Авторська рекомендація «У клініці» має підтверджений прийом.}{Автор пробує змінити формат на «Вдома», потім скасовує рекомендацію з причиною.}{Зміна формату відхилена, скасування збережене в історії; прийом не змінений. Інший лікар не редагує рекомендацію; подальше скасування прийому не повертає скасовану рекомендацію в список.}

\scenario{BDD-046}{Строк рекомендації та запис в інший день}{Рекомендація пов’язана з прийомом 28.09.2026, строк становить 10 календарних днів.}{Лікар зберігає рекомендацію, клієнт відкриває форму й потім успішно записується на інший доступний день.}{Рекомендована дата 08.10.2026 і строк 10 збережені. В окремих прогонах строк 1 день дає 29.09.2026, а значення 0, -1 і 1,5 відхиляються; верхня межа для додатного цілого строку не застосовується. Відкриття форми не прибирає рекомендацію зі списку, успішний запис прибирає. Без базового прийому чи процедури відносний строк не приймається. В окремому прогоні за прийому 01.10.2026 і процедури 03.10.2026 лікар явно обирає процедуру та строк 10 днів: дата рекомендації — 13.10.2026. Після виправлення дати процедури на 04.10.2026 рекомендація залишається на 13.10.2026; її може явно виправити автор із причиною та попередньою версією.}

\scenario{BDD-047}{Домашні та неактивні рекомендації}{Є домашня рекомендація з інструкцією й датою, скасована клінічна рекомендація та чинна клінічна для померлої тварини.}{Клієнт переглядає рекомендації, адміністратор відкриває список очікування.}{Домашня доступна для читання без онлайн-запису й позначки виконання. Жодна з трьох не показується у списку очікування; історія зберігається.}

\scenario{BDD-048}{Одна рекомендація та два різні часи}{За чинною рекомендацією ще немає підтвердженого або проведеного прийому.}{Клієнт і адміністратор одночасно створюють за нею записи на два різні вільні часи.}{Збережено лише один підтверджений прийом. Після його скасування або неявки повтор дозволений, а після проведення — ні.}

\scenario{BDD-049}{Непідтверджене надсилання SMS}{Користувач запитує код для однієї з трьох погоджених дій.}{SMS-сервіс відхиляє запит або не відповідає 10 секунд.}{Показано помилку або непідтверджене надсилання; дія не завершена, дані доступу не змінені. Повтор доступний за правилом 60 секунд. Прийняття запиту в іншому прогоні не вважається доставкою чи перевіркою номера.}

\scenario{BDD-050}{Помилка копіювання та пробне відновлення}{Налаштовані окреме резервне сховище й адреса технічної особи.}{Виявлено помилку копіювання, а в окремій контрольній вправі модельовано серйозний збій і відновлення.}{Про помилку надіслано лист до 5 хвилин від виявлення. Вправа відновлює узгоджені картки й файли за 4 години від зупинки з втратою змін не більше години; результати документовані.}

\scenario{BDD-051}{Бездіяльність під час введення тексту}{Працівник відкрив медичний запис і вводить текст без надсилання запитів; фонові оновлення тривають.}{Минає 30 хвилин від останньої явної дії із запитом, після чого він пробує зберегти.}{Сеанс завершений, потрібний новий вхід; текст не збережений автоматично. Окремий прогін для клієнта застосовує 24 години, а явний пошук до межі продовжує строк.}

\section{Початкова матриця простежуваності}
Кожна вимога має один стабільний AC із набором спостережуваних перевірок. BDD наведено там, де він змістовний; для процесних вимог, огляду конфігурації та вимірювань застосовуються безпосередні AC. Для цієї лабораторної роботи матриця містить зв’язки «вимога → критерій прийняття → BDD-сценарій», якщо сценарій є. Зв’язок тестів із вимогами забезпечується в межах TDD. Обов’язкове майбутнє доповнення матриці посиланнями на тікети, код і документацію не встановлюється.

\small
\begin{longtable}{p{3.1cm} p{3.0cm} >{\raggedright\arraybackslash}p{9.1cm}}
\toprule \textbf{Вимога} & \textbf{Критерій прийняття} & \textbf{BDD сценарії} \\ \midrule
\endfirsthead
\toprule \textbf{Вимога} & \textbf{Критерій прийняття} & \textbf{BDD сценарії} \\ \midrule
\endhead \bottomrule \endfoot
\hyperlink{REQ-F-001}{\texttt{REQ-F-001}} & \hyperlink{AC-F-001}{\texttt{AC-F-001}} & \hyperlink{BDD-001}{\texttt{BDD-001}}, \hyperlink{BDD-018}{\texttt{BDD-018}} \\
\hyperlink{REQ-F-002}{\texttt{REQ-F-002}} & \hyperlink{AC-F-002}{\texttt{AC-F-002}} & \hyperlink{BDD-001}{\texttt{BDD-001}}, \hyperlink{BDD-002}{\texttt{BDD-002}} \\
\hyperlink{REQ-F-003}{\texttt{REQ-F-003}} & \hyperlink{AC-F-003}{\texttt{AC-F-003}} & \hyperlink{BDD-001}{\texttt{BDD-001}}, \hyperlink{BDD-002}{\texttt{BDD-002}}, \hyperlink{BDD-003}{\texttt{BDD-003}} \\
\hyperlink{REQ-F-004}{\texttt{REQ-F-004}} & \hyperlink{AC-F-004}{\texttt{AC-F-004}} & \hyperlink{BDD-004}{\texttt{BDD-004}} \\
\hyperlink{REQ-F-005}{\texttt{REQ-F-005}} & \hyperlink{AC-F-005}{\texttt{AC-F-005}} & \hyperlink{BDD-005}{\texttt{BDD-005}} \\
\hyperlink{REQ-F-006}{\texttt{REQ-F-006}} & \hyperlink{AC-F-006}{\texttt{AC-F-006}} & \hyperlink{BDD-003}{\texttt{BDD-003}}, \hyperlink{BDD-006}{\texttt{BDD-006}} \\
\hyperlink{REQ-F-007}{\texttt{REQ-F-007}} & \hyperlink{AC-F-007}{\texttt{AC-F-007}} & \hyperlink{BDD-006}{\texttt{BDD-006}} \\
\hyperlink{REQ-F-008}{\texttt{REQ-F-008}} & \hyperlink{AC-F-008}{\texttt{AC-F-008}} & \hyperlink{BDD-007}{\texttt{BDD-007}} \\
\hyperlink{REQ-F-009}{\texttt{REQ-F-009}} & \hyperlink{AC-F-009}{\texttt{AC-F-009}} & \hyperlink{BDD-008}{\texttt{BDD-008}} \\
\hyperlink{REQ-F-010}{\texttt{REQ-F-010}} & \hyperlink{AC-F-010}{\texttt{AC-F-010}} & \hyperlink{BDD-006}{\texttt{BDD-006}}, \hyperlink{BDD-007}{\texttt{BDD-007}}, \hyperlink{BDD-008}{\texttt{BDD-008}}, \hyperlink{BDD-009}{\texttt{BDD-009}} \\
\hyperlink{REQ-F-011}{\texttt{REQ-F-011}} & \hyperlink{AC-F-011}{\texttt{AC-F-011}} & — \\
\hyperlink{REQ-F-012}{\texttt{REQ-F-012}} & \hyperlink{AC-F-012}{\texttt{AC-F-012}} & \hyperlink{BDD-009}{\texttt{BDD-009}} \\
\hyperlink{REQ-F-013}{\texttt{REQ-F-013}} & \hyperlink{AC-F-013}{\texttt{AC-F-013}} & \hyperlink{BDD-009}{\texttt{BDD-009}} \\
\hyperlink{REQ-F-014}{\texttt{REQ-F-014}} & \hyperlink{AC-F-014}{\texttt{AC-F-014}} & \hyperlink{BDD-009}{\texttt{BDD-009}} \\
\hyperlink{REQ-F-015}{\texttt{REQ-F-015}} & \hyperlink{AC-F-015}{\texttt{AC-F-015}} & \hyperlink{BDD-010}{\texttt{BDD-010}} \\
\hyperlink{REQ-F-016}{\texttt{REQ-F-016}} & \hyperlink{AC-F-016}{\texttt{AC-F-016}} & \hyperlink{BDD-010}{\texttt{BDD-010}} \\
\hyperlink{REQ-F-017}{\texttt{REQ-F-017}} & \hyperlink{AC-F-017}{\texttt{AC-F-017}} & \hyperlink{BDD-011}{\texttt{BDD-011}} \\
\hyperlink{REQ-F-018}{\texttt{REQ-F-018}} & \hyperlink{AC-F-018}{\texttt{AC-F-018}} & \hyperlink{BDD-009}{\texttt{BDD-009}} \\
\hyperlink{REQ-F-019}{\texttt{REQ-F-019}} & \hyperlink{AC-F-019}{\texttt{AC-F-019}} & \hyperlink{BDD-012}{\texttt{BDD-012}} \\
\hyperlink{REQ-F-020}{\texttt{REQ-F-020}} & \hyperlink{AC-F-020}{\texttt{AC-F-020}} & \hyperlink{BDD-013}{\texttt{BDD-013}} \\
\hyperlink{REQ-F-021}{\texttt{REQ-F-021}} & \hyperlink{AC-F-021}{\texttt{AC-F-021}} & \hyperlink{BDD-001}{\texttt{BDD-001}}, \hyperlink{BDD-018}{\texttt{BDD-018}} \\
\hyperlink{REQ-F-022}{\texttt{REQ-F-022}} & \hyperlink{AC-F-022}{\texttt{AC-F-022}} & \hyperlink{BDD-012}{\texttt{BDD-012}}, \hyperlink{BDD-018}{\texttt{BDD-018}} \\
\hyperlink{REQ-F-023}{\texttt{REQ-F-023}} & \hyperlink{AC-F-023}{\texttt{AC-F-023}} & \hyperlink{BDD-014}{\texttt{BDD-014}} \\
\hyperlink{REQ-F-024}{\texttt{REQ-F-024}} & \hyperlink{AC-F-024}{\texttt{AC-F-024}} & \hyperlink{BDD-014}{\texttt{BDD-014}} \\
\hyperlink{REQ-F-025}{\texttt{REQ-F-025}} & \hyperlink{AC-F-025}{\texttt{AC-F-025}} & \hyperlink{BDD-012}{\texttt{BDD-012}}, \hyperlink{BDD-014}{\texttt{BDD-014}} \\
\hyperlink{REQ-F-026}{\texttt{REQ-F-026}} & \hyperlink{AC-F-026}{\texttt{AC-F-026}} & \hyperlink{BDD-015}{\texttt{BDD-015}} \\
\hyperlink{REQ-F-027}{\texttt{REQ-F-027}} & \hyperlink{AC-F-027}{\texttt{AC-F-027}} & \hyperlink{BDD-015}{\texttt{BDD-015}}, \hyperlink{BDD-013}{\texttt{BDD-013}} \\
\hyperlink{REQ-F-028}{\texttt{REQ-F-028}} & \hyperlink{AC-F-028}{\texttt{AC-F-028}} & \hyperlink{BDD-016}{\texttt{BDD-016}} \\
\hyperlink{REQ-F-029}{\texttt{REQ-F-029}} & \hyperlink{AC-F-029}{\texttt{AC-F-029}} & \hyperlink{BDD-017}{\texttt{BDD-017}} \\
\hyperlink{REQ-F-030}{\texttt{REQ-F-030}} & \hyperlink{AC-F-030}{\texttt{AC-F-030}} & \hyperlink{BDD-018}{\texttt{BDD-018}}, \hyperlink{BDD-020}{\texttt{BDD-020}} \\
\hyperlink{REQ-F-031}{\texttt{REQ-F-031}} & \hyperlink{AC-F-031}{\texttt{AC-F-031}} & \hyperlink{BDD-018}{\texttt{BDD-018}} \\
\hyperlink{REQ-F-032}{\texttt{REQ-F-032}} & \hyperlink{AC-F-032}{\texttt{AC-F-032}} & \hyperlink{BDD-018}{\texttt{BDD-018}}, \hyperlink{BDD-019}{\texttt{BDD-019}} \\
\hyperlink{REQ-F-033}{\texttt{REQ-F-033}} & \hyperlink{AC-F-033}{\texttt{AC-F-033}} & \hyperlink{BDD-020}{\texttt{BDD-020}} \\
\hyperlink{REQ-F-034}{\texttt{REQ-F-034}} & \hyperlink{AC-F-034}{\texttt{AC-F-034}} & \hyperlink{BDD-021}{\texttt{BDD-021}} \\
\hyperlink{REQ-F-035}{\texttt{REQ-F-035}} & \hyperlink{AC-F-035}{\texttt{AC-F-035}} & \hyperlink{BDD-021}{\texttt{BDD-021}} \\
\hyperlink{REQ-F-036}{\texttt{REQ-F-036}} & \hyperlink{AC-F-036}{\texttt{AC-F-036}} & \hyperlink{BDD-022}{\texttt{BDD-022}} \\
\hyperlink{REQ-F-037}{\texttt{REQ-F-037}} & \hyperlink{AC-F-037}{\texttt{AC-F-037}} & \hyperlink{BDD-022}{\texttt{BDD-022}}, \hyperlink{BDD-023}{\texttt{BDD-023}} \\
\hyperlink{REQ-F-038}{\texttt{REQ-F-038}} & \hyperlink{AC-F-038}{\texttt{AC-F-038}} & \hyperlink{BDD-023}{\texttt{BDD-023}} \\
\hyperlink{REQ-F-039}{\texttt{REQ-F-039}} & \hyperlink{AC-F-039}{\texttt{AC-F-039}} & \hyperlink{BDD-024}{\texttt{BDD-024}} \\
\hyperlink{REQ-F-040}{\texttt{REQ-F-040}} & \hyperlink{AC-F-040}{\texttt{AC-F-040}} & \hyperlink{BDD-024}{\texttt{BDD-024}} \\
\hyperlink{REQ-F-041}{\texttt{REQ-F-041}} & \hyperlink{AC-F-041}{\texttt{AC-F-041}} & \hyperlink{BDD-025}{\texttt{BDD-025}}, \hyperlink{BDD-026}{\texttt{BDD-026}}, \hyperlink{BDD-032}{\texttt{BDD-032}} \\
\hyperlink{REQ-F-042}{\texttt{REQ-F-042}} & \hyperlink{AC-F-042}{\texttt{AC-F-042}} & \hyperlink{BDD-027}{\texttt{BDD-027}} \\
\hyperlink{REQ-F-043}{\texttt{REQ-F-043}} & \hyperlink{AC-F-043}{\texttt{AC-F-043}} & \hyperlink{BDD-028}{\texttt{BDD-028}}, \hyperlink{BDD-029}{\texttt{BDD-029}} \\
\hyperlink{REQ-F-044}{\texttt{REQ-F-044}} & \hyperlink{AC-F-044}{\texttt{AC-F-044}} & \hyperlink{BDD-010}{\texttt{BDD-010}}, \hyperlink{BDD-029}{\texttt{BDD-029}} \\
\hyperlink{REQ-F-045}{\texttt{REQ-F-045}} & \hyperlink{AC-F-045}{\texttt{AC-F-045}} & \hyperlink{BDD-030}{\texttt{BDD-030}} \\
\hyperlink{REQ-F-046}{\texttt{REQ-F-046}} & \hyperlink{AC-F-046}{\texttt{AC-F-046}} & \hyperlink{BDD-031}{\texttt{BDD-031}} \\
\hyperlink{REQ-F-047}{\texttt{REQ-F-047}} & \hyperlink{AC-F-047}{\texttt{AC-F-047}} & \hyperlink{BDD-032}{\texttt{BDD-032}} \\
\hyperlink{REQ-F-048}{\texttt{REQ-F-048}} & \hyperlink{AC-F-048}{\texttt{AC-F-048}} & \hyperlink{BDD-030}{\texttt{BDD-030}} \\
\hyperlink{REQ-F-049}{\texttt{REQ-F-049}} & \hyperlink{AC-F-049}{\texttt{AC-F-049}} & \hyperlink{BDD-033}{\texttt{BDD-033}}, \hyperlink{BDD-034}{\texttt{BDD-034}} \\
\hyperlink{REQ-F-050}{\texttt{REQ-F-050}} & \hyperlink{AC-F-050}{\texttt{AC-F-050}} & \hyperlink{BDD-033}{\texttt{BDD-033}}, \hyperlink{BDD-034}{\texttt{BDD-034}} \\
\hyperlink{REQ-F-051}{\texttt{REQ-F-051}} & \hyperlink{AC-F-051}{\texttt{AC-F-051}} & \hyperlink{BDD-035}{\texttt{BDD-035}} \\
\hyperlink{REQ-F-052}{\texttt{REQ-F-052}} & \hyperlink{AC-F-052}{\texttt{AC-F-052}} & \hyperlink{BDD-036}{\texttt{BDD-036}} \\
\hyperlink{REQ-F-053}{\texttt{REQ-F-053}} & \hyperlink{AC-F-053}{\texttt{AC-F-053}} & \hyperlink{BDD-035}{\texttt{BDD-035}} \\
\hyperlink{REQ-F-054}{\texttt{REQ-F-054}} & \hyperlink{AC-F-054}{\texttt{AC-F-054}} & \hyperlink{BDD-037}{\texttt{BDD-037}}, \hyperlink{BDD-038}{\texttt{BDD-038}} \\
\hyperlink{REQ-F-055}{\texttt{REQ-F-055}} & \hyperlink{AC-F-055}{\texttt{AC-F-055}} & \hyperlink{BDD-031}{\texttt{BDD-031}} \\
\hyperlink{REQ-F-056}{\texttt{REQ-F-056}} & \hyperlink{AC-F-056}{\texttt{AC-F-056}} & \hyperlink{BDD-039}{\texttt{BDD-039}} \\
\hyperlink{REQ-F-057}{\texttt{REQ-F-057}} & \hyperlink{AC-F-057}{\texttt{AC-F-057}} & \hyperlink{BDD-040}{\texttt{BDD-040}} \\
\hyperlink{REQ-F-058}{\texttt{REQ-F-058}} & \hyperlink{AC-F-058}{\texttt{AC-F-058}} & \hyperlink{BDD-040}{\texttt{BDD-040}}, \hyperlink{BDD-041}{\texttt{BDD-041}} \\
\hyperlink{REQ-F-059}{\texttt{REQ-F-059}} & \hyperlink{AC-F-059}{\texttt{AC-F-059}} & \hyperlink{BDD-041}{\texttt{BDD-041}}, \hyperlink{BDD-042}{\texttt{BDD-042}} \\
\hyperlink{REQ-F-060}{\texttt{REQ-F-060}} & \hyperlink{AC-F-060}{\texttt{AC-F-060}} & \hyperlink{BDD-042}{\texttt{BDD-042}} \\
\hyperlink{REQ-F-061}{\texttt{REQ-F-061}} & \hyperlink{AC-F-061}{\texttt{AC-F-061}} & \hyperlink{BDD-043}{\texttt{BDD-043}} \\
\hyperlink{REQ-F-062}{\texttt{REQ-F-062}} & \hyperlink{AC-F-062}{\texttt{AC-F-062}} & \hyperlink{BDD-043}{\texttt{BDD-043}} \\
\hyperlink{REQ-F-063}{\texttt{REQ-F-063}} & \hyperlink{AC-F-063}{\texttt{AC-F-063}} & \hyperlink{BDD-037}{\texttt{BDD-037}} \\
\hyperlink{REQ-F-064}{\texttt{REQ-F-064}} & \hyperlink{AC-F-064}{\texttt{AC-F-064}} & \hyperlink{BDD-037}{\texttt{BDD-037}} \\
\hyperlink{REQ-F-065}{\texttt{REQ-F-065}} & \hyperlink{AC-F-065}{\texttt{AC-F-065}} & \hyperlink{BDD-044}{\texttt{BDD-044}} \\
\hyperlink{REQ-F-066}{\texttt{REQ-F-066}} & \hyperlink{AC-F-066}{\texttt{AC-F-066}} & \hyperlink{BDD-045}{\texttt{BDD-045}} \\
\hyperlink{REQ-F-067}{\texttt{REQ-F-067}} & \hyperlink{AC-F-067}{\texttt{AC-F-067}} & \hyperlink{BDD-046}{\texttt{BDD-046}} \\
\hyperlink{REQ-F-068}{\texttt{REQ-F-068}} & \hyperlink{AC-F-068}{\texttt{AC-F-068}} & \hyperlink{BDD-046}{\texttt{BDD-046}}, \hyperlink{BDD-047}{\texttt{BDD-047}} \\
\hyperlink{REQ-F-069}{\texttt{REQ-F-069}} & \hyperlink{AC-F-069}{\texttt{AC-F-069}} & \hyperlink{BDD-046}{\texttt{BDD-046}}, \hyperlink{BDD-048}{\texttt{BDD-048}} \\
\hyperlink{REQ-F-070}{\texttt{REQ-F-070}} & \hyperlink{AC-F-070}{\texttt{AC-F-070}} & \hyperlink{BDD-032}{\texttt{BDD-032}}, \hyperlink{BDD-034}{\texttt{BDD-034}}, \hyperlink{BDD-047}{\texttt{BDD-047}} \\
\hyperlink{REQ-F-071}{\texttt{REQ-F-071}} & \hyperlink{AC-F-071}{\texttt{AC-F-071}} & \hyperlink{BDD-045}{\texttt{BDD-045}}, \hyperlink{BDD-047}{\texttt{BDD-047}}, \hyperlink{BDD-028}{\texttt{BDD-028}} \\
\hyperlink{REQ-F-072}{\texttt{REQ-F-072}} & \hyperlink{AC-F-072}{\texttt{AC-F-072}} & \hyperlink{BDD-047}{\texttt{BDD-047}} \\
\hyperlink{REQ-F-073}{\texttt{REQ-F-073}} & \hyperlink{AC-F-073}{\texttt{AC-F-073}} & \hyperlink{BDD-049}{\texttt{BDD-049}} \\
\hyperlink{REQ-F-074}{\texttt{REQ-F-074}} & \hyperlink{AC-F-074}{\texttt{AC-F-074}} & \hyperlink{BDD-050}{\texttt{BDD-050}} \\
\hyperlink{REQ-NF-001}{\texttt{REQ-NF-001}} & \hyperlink{AC-NF-001}{\texttt{AC-NF-001}} & — \\
\hyperlink{REQ-NF-002}{\texttt{REQ-NF-002}} & \hyperlink{AC-NF-002}{\texttt{AC-NF-002}} & — \\
\hyperlink{REQ-NF-003}{\texttt{REQ-NF-003}} & \hyperlink{AC-NF-003}{\texttt{AC-NF-003}} & \hyperlink{BDD-022}{\texttt{BDD-022}}, \hyperlink{BDD-036}{\texttt{BDD-036}} \\
\hyperlink{REQ-NF-004}{\texttt{REQ-NF-004}} & \hyperlink{AC-NF-004}{\texttt{AC-NF-004}} & — \\
\hyperlink{REQ-NF-005}{\texttt{REQ-NF-005}} & \hyperlink{AC-NF-005}{\texttt{AC-NF-005}} & — \\
\hyperlink{REQ-NF-006}{\texttt{REQ-NF-006}} & \hyperlink{AC-NF-006}{\texttt{AC-NF-006}} & \hyperlink{BDD-049}{\texttt{BDD-049}} \\
\hyperlink{REQ-NF-007}{\texttt{REQ-NF-007}} & \hyperlink{AC-NF-007}{\texttt{AC-NF-007}} & \hyperlink{BDD-050}{\texttt{BDD-050}} \\
\hyperlink{REQ-NF-008}{\texttt{REQ-NF-008}} & \hyperlink{AC-NF-008}{\texttt{AC-NF-008}} & — \\
\hyperlink{REQ-NF-009}{\texttt{REQ-NF-009}} & \hyperlink{AC-NF-009}{\texttt{AC-NF-009}} & \hyperlink{BDD-050}{\texttt{BDD-050}} \\
\hyperlink{REQ-NF-010}{\texttt{REQ-NF-010}} & \hyperlink{AC-NF-010}{\texttt{AC-NF-010}} & — \\
\hyperlink{REQ-NF-011}{\texttt{REQ-NF-011}} & \hyperlink{AC-NF-011}{\texttt{AC-NF-011}} & \hyperlink{BDD-005}{\texttt{BDD-005}} \\
\hyperlink{REQ-NF-012}{\texttt{REQ-NF-012}} & \hyperlink{AC-NF-012}{\texttt{AC-NF-012}} & \hyperlink{BDD-051}{\texttt{BDD-051}} \\
\hyperlink{REQ-NF-013}{\texttt{REQ-NF-013}} & \hyperlink{AC-NF-013}{\texttt{AC-NF-013}} & — \\
\hyperlink{REQ-NF-014}{\texttt{REQ-NF-014}} & \hyperlink{AC-NF-014}{\texttt{AC-NF-014}} & — \\
\hyperlink{REQ-NF-015}{\texttt{REQ-NF-015}} & \hyperlink{AC-NF-015}{\texttt{AC-NF-015}} & \hyperlink{BDD-038}{\texttt{BDD-038}} \\
\hyperlink{REQ-DP-001}{\texttt{REQ-DP-001}} & \hyperlink{AC-DP-001}{\texttt{AC-DP-001}} & — \\
\hyperlink{REQ-DP-002}{\texttt{REQ-DP-002}} & \hyperlink{AC-DP-002}{\texttt{AC-DP-002}} & — \\
\hyperlink{REQ-DP-003}{\texttt{REQ-DP-003}} & \hyperlink{AC-DP-003}{\texttt{AC-DP-003}} & — \\
\hyperlink{REQ-DP-004}{\texttt{REQ-DP-004}} & \hyperlink{AC-DP-004}{\texttt{AC-DP-004}} & — \\
\hyperlink{REQ-DP-005}{\texttt{REQ-DP-005}} & \hyperlink{AC-DP-005}{\texttt{AC-DP-005}} & — \\
\hyperlink{REQ-DP-006}{\texttt{REQ-DP-006}} & \hyperlink{AC-DP-006}{\texttt{AC-DP-006}} & — \\
\end{longtable}
\normalsize
\section{Порядок перевірки та матеріали прийняття}
\begin{itemize}
\item Перевірити рольову матрицю прямими запитами, а не лише видимістю кнопок; застосувати негативні BDD до файлів і попередніх версій.
\item Виконати AC для всіх REQ-F і релевантні BDD. Часові тести використовують керований час, конкурентні — дійсно одночасні спроби, а не лише послідовні кліки.
\item Окремо перевірити межі кодів, посилань, блокування й сеансів, унікальність телефону та мікрочипа, інтервали й атомарність бронювання.
\item Для продуктивності перед прогоном зафіксувати DEP-001: середовище, мережу, розподіл дій, кількість спроб і метод вимірювання. Пороги оцінювати окремо, не середнім для всіх типів.
\item Виконати окремий інтеграційний тест доставки SMS на контрольний номер і надсилання технічного листа. Модельована відповідь сервісу не доводить фактичну доставку.
\item Провести пробне резервне відновлення в окремому середовищі, перевіривши час, втрату змін, картки й вкладення. Зберегти результати.
\item Перевірити виконання REQ-DP за доказами до і після реалізації. Не називати структурну перевірку LaTeX або матриці доказом реалізованої системи.
\item Окремо провести SPEC-GATE: перевірити OPEN і DEP, суперечності, терміни, критерії, межі та простежуваність. Після перевірки зафіксувати PASS або FAIL; за FAIL опрацювати блокуючі зауваження й повторити перевірку. Після PASS отримати людське затвердження перед комітом затвердженої baseline.
\end{itemize}

\section{Залежності та питання для SPEC-GATE}
Наступні записи є чесно позначеними залежностями й деталями, виявленими при формалізації. Погоджені після SPEC-GATE рішення зазначені у відповідних рядках зі збереженням їхніх ID; решта невизначеностей і залежностей залишаються відкритими. Перед затвердженням потрібно вирішити питання, що впливають на поведінку, або явно відкласти відповідний сценарій. Технічні залежності можна закрити до належного етапу реалізації чи приймання за окремим погодженням.

\begin{longtable}{>{\raggedright\arraybackslash}p{2.4cm} >{\raggedright\arraybackslash}p{3.2cm} >{\raggedright\arraybackslash}p{9.6cm}}
\toprule
\textbf{ID} & \textbf{Тема} & \textbf{Необхідна дія} \\ \midrule
\endfirsthead
\toprule
\textbf{ID} & \textbf{Тема} & \textbf{Необхідна дія} \\ \midrule
\endhead
\bottomrule
\endfoot
DEP-001 & Приймальне середовище & До вимірювання погодити сервер, сховище, мережу, точні ОС/браузери, розміри екранів, тривалість і профіль навантаження, розподіл розмірів карток/вкладень та кількість спроб. Числові пороги вже погоджені, методика ще не повністю задана. \\
DEP-002 & Постачальники інтеграцій & До інтеграційних тестів обрати SMS і технічну пошту, контрольний номер, адресу відповідальної особи й облікові дані сервісів. Секрети не включаються до SRS. \\
DEP-003 & Розгортання і резервування & До експлуатації визначити хостинг, окреме сховище, процедуру виявлення збоїв, відповідальну технічну особу та інструкцію відновлення. \\
DEP-004 & Перевірка особи & Клініка до реального використання визначає достатність доказів при особистому зверненні. Система фіксує рішення адміністратора, а не автоматично доводить особу за збігом імен. \\
DEP-005 & Рішення реалізації захисту & До реалізації визначити алгоритм і параметри хешування, зберігання секретів інтеграцій та конфігурацію HTTPS. Суттєві зміни поведінки мають повертатися на погодження. \\
\end{longtable}

\begin{longtable}{>{\raggedright\arraybackslash}p{2.4cm} >{\raggedright\arraybackslash}p{3.2cm} >{\raggedright\arraybackslash}p{9.6cm}}
\toprule
\textbf{ID} & \textbf{Тема} & \textbf{Невизначеність} \\ \midrule
\endfirsthead
\toprule
\textbf{ID} & \textbf{Тема} & \textbf{Невизначеність} \\ \midrule
\endhead
\bottomrule
\endfoot
OPEN-001 & Відносний строк рекомендації & Закрито за SG-B06: лікар явно обирає одну підставу; зміна її дати не перераховує збережену рекомендацію автоматично (REQ-F-067). За рішенням людини від 02.10.2026 відносний строк — ціле число календарних днів від 1 і більше; 0, від’ємні та дробові значення не допускаються, верхню межу не встановлено. \\
OPEN-002 & Історичні прийоми та розклад & Рішення погоджено за SG-B04: перевіряються лише майбутні підтверджені прийоми, достатньо повного вміщення без збігу з новою сіткою; історичні прийоми не блокують зміну (REQ-F-038). \\
OPEN-003 & Єдиний адміністратор утратив пароль & Відновлення працівника потребує іншого адміністратора. Аварійний спосіб для єдиного активного адміністратора не погоджений; не вводити прихований привілейований вхід. \\
OPEN-004 & Проміжні стани відновлення & Рішення погоджено за SG-B07: посилання діє 24 години, підтвердження номера — до завершення його строку; прив’язка до профілю, номера й операції, заміна операції новим посиланням та незмінність старих даних визначені REQ-F-009. \\
OPEN-005 & Права на зміну рекомендацій та бронювання & Рішення погоджено за SG-B05: перенесення обома ролями лише на майбутній час; клієнт також може змінити лікаря за REQ-F-044. Скасована рекомендація не блокує перенесення чинного прийому, зв’язок і скасування зберігаються (REQ-F-043, REQ-F-071). \\
OPEN-006 & Збій технічної пошти & Погоджено лист за 5 хвилин, але не визначено поведінку, якщо недоступний сам поштовий сервіс. Не додавати резервний канал або необмежені повтори без рішення. \\
OPEN-007 & Точні правила валідації & Рішення погоджено за SG-B11: кодові точки Unicode без автоматичної зміни пароля; народження та події не можуть бути майбутніми; неоднозначний/неіснуючий місцевий час не приймається (REQ-F-019, REQ-F-030, REQ-NF-003, REQ-NF-010). Планові прийоми й рекомендації можуть мати майбутні дати; порожні/невідомі дані не підмінюються вигаданими. \\
OPEN-008 & Граничні права історичного відновлення & Рішення погоджено за SG-B02 і SG-B03: для тварини історичне створення враховує підтверджені та проведені прийоми незалежно від лікаря (REQ-F-041). Майбутнє відновлення застосовує обмеження майбутнього запису, минуле — історичні правила (REQ-F-050). \\
\end{longtable}

\section{Підсумок SKED і рішення щодо припущень}
Задум ЛР1 збережено: облік клієнтів і тварин, медична історія, планування та онлайн-запис, три ролі. Питання реєстрації до першої тварини закрито явним рішенням людини. Подальший діалог послідовно уточнив права, статуси, винятки, перевірки, інтеграції та процесні вимоги. Таблиця нижче фіксує приклади відхилених або виправлених пропозицій моделі, а не альтернативні чинні вимоги.

\begin{longtable}{>{\raggedright\arraybackslash}p{7.6cm} >{\raggedright\arraybackslash}p{7.6cm}}
\toprule
\textbf{Пропозиція чи початкова невизначеність} & \textbf{Рішення людини} \\ \midrule
\endfirsthead
\toprule
\textbf{Пропозиція чи початкова невизначеність} & \textbf{Рішення людини} \\ \midrule
\endhead
\bottomrule
\endfoot
Адміністратору закрити всю медичну інформацію & Дозволено організаційну частину клінічних рекомендацій; медичний зміст лишився недоступним. \\
Клієнту заборонити зміни після підтвердження тварини & Клієнт подає пропозицію; працівник приймає або відхиляє її цілком. \\
Окремий стан рекомендації «Виконана» & Не вводиться; враховується статус пов’язаного прийому. Скасування рекомендації погоджене окремо. \\
Один обліковий запис із перемиканням ролей & Обрано окремі робочі та особисті облікові записи. \\
Лише майбутній час для всіх ролей & Адміністратору дозволено історичне внесення поза розкладом без межі давності, з перевірками. \\
Заборонити історичний прийом заблокованого лікаря & Дозволено створення; прийом лишається підтвердженим до можливого розблокування призначеного лікаря. \\
Передавання картки іншому власнику & Не включено; автоматичного дублювання картки також немає. \\
Логи як загальне журналювання бізнес-подій & Логи викладача трактуються як процес розробки й докази TDD. Історії виправлень погоджені окремими функціями. \\
Незмінність кінцевих статусів & Погоджено контрольовані виправлення з правами, причинами й історією. \\
Редагування запису заблокованого автора іншим лікарем & Не дозволено; спеціального коригувального механізму у версії немає. \\
\end{longtable}

\section{Статус наступного кроку}
Ця поставка є затвердженою baseline версії 0.1 для подальшого SDD-циклу. SPEC-GATE від 02.10.2026 має результат PASS; усі блокуючі зауваження SG-B01–SG-B11 закриті. Абраменко Яна явно затвердила поточну специфікацію 03.10.2026. Коміт затвердженої базової специфікації ще не створено.

\end{document}
